%%%%%%%%%%%%%%%%%%%%%%%%%%%%%%%%%%%%%%%%%%%%%%%%%%%%%%%%%%%%
% 7.6 EXPECTATION AND MOMENTS
% The Composition of Advanced Mathematics
%%%%%%%%%%%%%%%%%%%%%%%%%%%%%%%%%%%%%%%%%%%%%%%%%%%%%%%%%%%%

\section{Expectation and Moments}
\label{sec:expectation-moments}

\prerequisite{
Random variables, discrete and continuous distributions, joint
distributions, probability mass functions, probability density
functions, infinite series, definite and improper integrals, and basic
algebra.
}

\learningobjectives{
By the end of this section, the reader should be able to:
\begin{itemize}
    \item define expected value for discrete and continuous random
          variables;
    \item interpret expectation as a probability-weighted average;
    \item determine whether an expectation exists;
    \item compute expectations using PMFs and PDFs;
    \item apply the law of the unconscious statistician;
    \item use linearity of expectation;
    \item compute expectations using indicator variables;
    \item define raw moments and central moments;
    \item define variance and standard deviation;
    \item use equivalent formulas for variance;
    \item understand the effects of linear transformations on mean and
          variance;
    \item compute means and variances of standard probability
          distributions;
    \item define covariance;
    \item derive covariance identities;
    \item define correlation;
    \item interpret the sign and magnitude of correlation;
    \item distinguish independence, uncorrelatedness, and perfect linear
          dependence;
    \item compute variance of sums of random variables;
    \item understand covariance matrices;
    \item define conditional expectation;
    \item use the law of total expectation;
    \item use the law of total variance;
    \item define higher moments, skewness, and kurtosis;
    \item understand factorial moments;
    \item recognize situations where moments fail to exist;
    \item apply moment inequalities such as Markov and Chebyshev;
    \item connect moments with transforms, limit theorems, statistics,
          and stochastic processes.
\end{itemize}
}

%%%%%%%%%%%%%%%%%%%%%%%%%%%%%%%%%%%%%%%%%%%%%%%%%%%%%%%%%%%%
\subsection{What Is Expected Value?}
\label{subsec:what-is-expected-value}
%%%%%%%%%%%%%%%%%%%%%%%%%%%%%%%%%%%%%%%%%%%%%%%%%%%%%%%%%%%%

The expected value of a random variable is a probability-weighted
average of its possible values.

It is often interpreted as the long-run average value that would emerge
if the same probabilistic experiment were repeated many times.

The expected value of \(X\) is written

\[
\boxed{
\mathbb{E}[X].
}
\]

The symbol

\[
\mathbb{E}
\]

is read ``expectation.''

One also commonly writes

\[
\boxed{
\mu_X
=
\mathbb{E}[X].
}
\]

%%%%%%%%%%%%%%%%%%%%%%%%%%%%%%%%%%%%%%%%%%%%%%%%%%%%%%%%%%%%
\subsection{Expected Value of a Discrete Random Variable}
\label{subsec:discrete-expectation}
%%%%%%%%%%%%%%%%%%%%%%%%%%%%%%%%%%%%%%%%%%%%%%%%%%%%%%%%%%%%

\begin{definitionbox}{Expected Value: Discrete Case}
If \(X\) is discrete with PMF

\[
p_X(x)=P(X=x),
\]

then

\[
\boxed{
\mathbb{E}[X]
=
\sum_x
x\,p_X(x),
}
\]

provided the sum is well-defined.
\end{definitionbox}

Each possible value is multiplied by its probability.

%%%%%%%%%%%%%%%%%%%%%%%%%%%%%%%%%%%%%%%%%%%%%%%%%%%%%%%%%%%%
\subsection{Worked Example: Expected Die Result}
\label{subsec:worked-expected-die}
%%%%%%%%%%%%%%%%%%%%%%%%%%%%%%%%%%%%%%%%%%%%%%%%%%%%%%%%%%%%

\begin{workedexamplebox}{Mean of a Fair Die}

Let \(X\) be the result of a fair six-sided die.

Then

\[
P(X=x)=\frac16,
\qquad
x=1,\ldots,6.
\]

Therefore,

\[
\mathbb{E}[X]
=
\sum_{x=1}^{6}
x\left(\frac16\right).
\]

Thus,

\[
\mathbb{E}[X]
=
\frac{1+2+3+4+5+6}{6}.
\]

\begin{answerbox}
\[
\boxed{
\mathbb{E}[X]
=
\frac72.
}
\]
\end{answerbox}

\end{workedexamplebox}

%%%%%%%%%%%%%%%%%%%%%%%%%%%%%%%%%%%%%%%%%%%%%%%%%%%%%%%%%%%%
\subsection{Expectation Need Not Be a Possible Value}
\label{subsec:expectation-not-possible-value}
%%%%%%%%%%%%%%%%%%%%%%%%%%%%%%%%%%%%%%%%%%%%%%%%%%%%%%%%%%%%

The expected value does not need to belong to the support.

For a fair die,

\[
\mathbb{E}[X]
=
3.5,
\]

even though no roll can equal \(3.5\).

Thus,

\[
\boxed{
\text{expected value}
\neq
\text{necessarily an attainable outcome}.
}
\]

It is a weighted center of the distribution.

%%%%%%%%%%%%%%%%%%%%%%%%%%%%%%%%%%%%%%%%%%%%%%%%%%%%%%%%%%%%
\subsection{Expected Value of a Continuous Random Variable}
\label{subsec:continuous-expectation}
%%%%%%%%%%%%%%%%%%%%%%%%%%%%%%%%%%%%%%%%%%%%%%%%%%%%%%%%%%%%

\begin{definitionbox}{Expected Value: Continuous Case}
If \(X\) has probability density \(f_X\), then

\[
\boxed{
\mathbb{E}[X]
=
\int_{-\infty}^{\infty}
x f_X(x)\,dx,
}
\]

provided the integral exists.
\end{definitionbox}

The integral is the continuous analogue of the weighted sum used in the
discrete case.

%%%%%%%%%%%%%%%%%%%%%%%%%%%%%%%%%%%%%%%%%%%%%%%%%%%%%%%%%%%%
\subsection{Worked Example: Uniform Mean}
\label{subsec:worked-uniform-mean}
%%%%%%%%%%%%%%%%%%%%%%%%%%%%%%%%%%%%%%%%%%%%%%%%%%%%%%%%%%%%

\begin{workedexamplebox}{Mean of a Uniform Distribution}

Suppose

\[
X\sim\operatorname{Uniform}(a,b).
\]

Then

\[
f_X(x)
=
\frac1{b-a}
\]

for

\[
a\leq x\leq b.
\]

Therefore,

\[
\mathbb{E}[X]
=
\int_a^b
x\frac1{b-a}\,dx.
\]

Thus,

\[
=
\frac1{b-a}
\left[
\frac{x^2}{2}
\right]_a^b.
\]

Hence,

\begin{answerbox}
\[
\boxed{
\mathbb{E}[X]
=
\frac{a+b}{2}.
}
\]
\end{answerbox}

\end{workedexamplebox}

%%%%%%%%%%%%%%%%%%%%%%%%%%%%%%%%%%%%%%%%%%%%%%%%%%%%%%%%%%%%
\subsection{Existence of Expectation}
\label{subsec:expectation-existence}
%%%%%%%%%%%%%%%%%%%%%%%%%%%%%%%%%%%%%%%%%%%%%%%%%%%%%%%%%%%%

Expectation is not automatically finite.

A useful sufficient condition is

\[
\boxed{
\mathbb{E}[|X|]<\infty.
}
\]

For a discrete variable,

\[
\mathbb{E}[|X|]
=
\sum_x
|x|p_X(x).
\]

For a continuous variable,

\[
\mathbb{E}[|X|]
=
\int_{-\infty}^{\infty}
|x|f_X(x)\,dx.
\]

If this quantity is finite, then \(\mathbb{E}[X]\) exists as a finite
real number.

%%%%%%%%%%%%%%%%%%%%%%%%%%%%%%%%%%%%%%%%%%%%%%%%%%%%%%%%%%%%
\subsection{Expectation of a Function}
\label{subsec:expectation-function-rv}
%%%%%%%%%%%%%%%%%%%%%%%%%%%%%%%%%%%%%%%%%%%%%%%%%%%%%%%%%%%%

Often one wants

\[
\mathbb{E}[g(X)]
\]

for some function \(g\).

It is not necessary to first find the distribution of \(g(X)\).

%%%%%%%%%%%%%%%%%%%%%%%%%%%%%%%%%%%%%%%%%%%%%%%%%%%%%%%%%%%%
\subsection{Law of the Unconscious Statistician}
\label{subsec:lotus}
%%%%%%%%%%%%%%%%%%%%%%%%%%%%%%%%%%%%%%%%%%%%%%%%%%%%%%%%%%%%

\begin{theorem}[Law of the Unconscious Statistician]
If \(X\) is discrete, then

\[
\boxed{
\mathbb{E}[g(X)]
=
\sum_x
g(x)p_X(x).
}
\]

If \(X\) is continuous with density \(f_X\), then

\[
\boxed{
\mathbb{E}[g(X)]
=
\int_{-\infty}^{\infty}
g(x)f_X(x)\,dx.
}
\]
\end{theorem}

This result is often abbreviated informally as LOTUS.

%%%%%%%%%%%%%%%%%%%%%%%%%%%%%%%%%%%%%%%%%%%%%%%%%%%%%%%%%%%%
\subsection{Worked Example: Expected Square}
\label{subsec:worked-expected-square}
%%%%%%%%%%%%%%%%%%%%%%%%%%%%%%%%%%%%%%%%%%%%%%%%%%%%%%%%%%%%

\begin{workedexamplebox}{Square of a Fair Die}

Let \(X\) be a fair die result.

Then

\[
\mathbb{E}[X^2]
=
\sum_{x=1}^{6}
x^2\frac16.
\]

Therefore,

\[
\mathbb{E}[X^2]
=
\frac{
1^2+2^2+3^2+4^2+5^2+6^2
}{6}.
\]

Since

\[
1+4+9+16+25+36=91,
\]

\begin{answerbox}
\[
\boxed{
\mathbb{E}[X^2]
=
\frac{91}{6}.
}
\]
\end{answerbox}

\end{workedexamplebox}

%%%%%%%%%%%%%%%%%%%%%%%%%%%%%%%%%%%%%%%%%%%%%%%%%%%%%%%%%%%%
\subsection{Linearity of Expectation}
\label{subsec:linearity-expectation}
%%%%%%%%%%%%%%%%%%%%%%%%%%%%%%%%%%%%%%%%%%%%%%%%%%%%%%%%%%%%

One of the most useful facts in probability is linearity of
expectation.

\begin{theorem}[Linearity of Expectation]
For random variables \(X_1,\ldots,X_n\) with finite expectations and
constants \(a_1,\ldots,a_n,b\),

\[
\boxed{
\mathbb{E}
\left[
b+\sum_{i=1}^{n}a_iX_i
\right]
=
b+\sum_{i=1}^{n}a_i\mathbb{E}[X_i].
}
\]
\end{theorem}

Independence is not required.

%%%%%%%%%%%%%%%%%%%%%%%%%%%%%%%%%%%%%%%%%%%%%%%%%%%%%%%%%%%%
\subsection{Two-Variable Linearity}
\label{subsec:two-variable-linearity}
%%%%%%%%%%%%%%%%%%%%%%%%%%%%%%%%%%%%%%%%%%%%%%%%%%%%%%%%%%%%

In particular,

\[
\boxed{
\mathbb{E}[X+Y]
=
\mathbb{E}[X]
+
\mathbb{E}[Y].
}
\]

Also,

\[
\boxed{
\mathbb{E}[aX+b]
=
a\mathbb{E}[X]+b.
}
\]

These formulas hold regardless of dependence.

%%%%%%%%%%%%%%%%%%%%%%%%%%%%%%%%%%%%%%%%%%%%%%%%%%%%%%%%%%%%
\subsection{Worked Example: Total Cost}
\label{subsec:worked-linearity-cost}
%%%%%%%%%%%%%%%%%%%%%%%%%%%%%%%%%%%%%%%%%%%%%%%%%%%%%%%%%%%%

\begin{workedexamplebox}{Linear Cost Transformation}

Suppose \(X\) is a random number of units used and

\[
\mathbb{E}[X]=40.
\]

A cost is defined by

\[
C=5X+12.
\]

Then

\[
\mathbb{E}[C]
=
5\mathbb{E}[X]+12.
\]

Therefore,

\begin{answerbox}
\[
\boxed{
\mathbb{E}[C]
=
212.
}
\]
\end{answerbox}

\end{workedexamplebox}

%%%%%%%%%%%%%%%%%%%%%%%%%%%%%%%%%%%%%%%%%%%%%%%%%%%%%%%%%%%%
\subsection{Indicators and Expectation}
\label{subsec:indicator-expectation}
%%%%%%%%%%%%%%%%%%%%%%%%%%%%%%%%%%%%%%%%%%%%%%%%%%%%%%%%%%%%

For an event \(A\),

\[
\mathbf{1}_A
=
\begin{cases}
1,&A\text{ occurs},\\
0,&A\text{ does not occur}.
\end{cases}
\]

Then

\[
\mathbb{E}[\mathbf{1}_A]
=
1\cdot P(A)
+
0\cdot P(A^c).
\]

Therefore,

\[
\boxed{
\mathbb{E}[\mathbf{1}_A]
=
P(A).
}
\]

This is one of the most useful identities in combinatorial probability.

%%%%%%%%%%%%%%%%%%%%%%%%%%%%%%%%%%%%%%%%%%%%%%%%%%%%%%%%%%%%
\subsection{Counting with Indicators}
\label{subsec:counting-with-indicators}
%%%%%%%%%%%%%%%%%%%%%%%%%%%%%%%%%%%%%%%%%%%%%%%%%%%%%%%%%%%%

Suppose

\[
X
=
\sum_{i=1}^{n}
\mathbf{1}_{A_i}.
\]

Then \(X\) counts the number of events \(A_i\) that occur.

By linearity,

\[
\boxed{
\mathbb{E}[X]
=
\sum_{i=1}^{n}
P(A_i).
}
\]

No independence is required.

%%%%%%%%%%%%%%%%%%%%%%%%%%%%%%%%%%%%%%%%%%%%%%%%%%%%%%%%%%%%
\subsection{Worked Example: Fixed Points of a Random Permutation}
\label{subsec:worked-fixed-points-expectation}
%%%%%%%%%%%%%%%%%%%%%%%%%%%%%%%%%%%%%%%%%%%%%%%%%%%%%%%%%%%%

\begin{workedexamplebox}{Expected Number of Fixed Points}

Choose a permutation of

\[
\{1,\ldots,n\}
\]

uniformly at random.

Let

\[
I_i
=
\mathbf{1}_{\{\pi(i)=i\}}.
\]

Then the number of fixed points is

\[
X
=
\sum_{i=1}^{n}I_i.
\]

For each \(i\),

\[
P(\pi(i)=i)
=
\frac1n.
\]

Therefore,

\[
\mathbb{E}[X]
=
\sum_{i=1}^{n}
\frac1n.
\]

\begin{answerbox}
\[
\boxed{
\mathbb{E}[X]=1.
}
\]
\end{answerbox}

\end{workedexamplebox}

%%%%%%%%%%%%%%%%%%%%%%%%%%%%%%%%%%%%%%%%%%%%%%%%%%%%%%%%%%%%
\subsection{Tail-Sum Formula}
\label{subsec:tail-sum-formula}
%%%%%%%%%%%%%%%%%%%%%%%%%%%%%%%%%%%%%%%%%%%%%%%%%%%%%%%%%%%%

For a nonnegative integer-valued random variable,

\[
\boxed{
\mathbb{E}[X]
=
\sum_{k=1}^{\infty}
P(X\geq k).
}
\]

Equivalently,

\[
\boxed{
\mathbb{E}[X]
=
\sum_{k=0}^{\infty}
P(X>k).
}
\]

This is called the tail-sum formula.

%%%%%%%%%%%%%%%%%%%%%%%%%%%%%%%%%%%%%%%%%%%%%%%%%%%%%%%%%%%%
\subsection{Continuous Tail Formula}
\label{subsec:continuous-tail-formula}
%%%%%%%%%%%%%%%%%%%%%%%%%%%%%%%%%%%%%%%%%%%%%%%%%%%%%%%%%%%%

For a nonnegative continuous random variable,

\[
\boxed{
\mathbb{E}[X]
=
\int_0^\infty
P(X>x)\,dx.
}
\]

Since

\[
P(X>x)=S_X(x),
\]

we may write

\[
\boxed{
\mathbb{E}[X]
=
\int_0^\infty
S_X(x)\,dx.
}
\]

%%%%%%%%%%%%%%%%%%%%%%%%%%%%%%%%%%%%%%%%%%%%%%%%%%%%%%%%%%%%
\subsection{Worked Example: Exponential Mean from the Tail}
\label{subsec:worked-exponential-tail-mean}
%%%%%%%%%%%%%%%%%%%%%%%%%%%%%%%%%%%%%%%%%%%%%%%%%%%%%%%%%%%%

\begin{workedexamplebox}{Mean Using the Survival Function}

Suppose

\[
X\sim\operatorname{Exp}(\lambda).
\]

Then

\[
P(X>x)
=
e^{-\lambda x}.
\]

Therefore,

\[
\mathbb{E}[X]
=
\int_0^\infty
e^{-\lambda x}\,dx.
\]

Thus,

\begin{answerbox}
\[
\boxed{
\mathbb{E}[X]
=
\frac1\lambda.
}
\]
\end{answerbox}

\end{workedexamplebox}

%%%%%%%%%%%%%%%%%%%%%%%%%%%%%%%%%%%%%%%%%%%%%%%%%%%%%%%%%%%%
\subsection{Moments}
\label{subsec:moments}
%%%%%%%%%%%%%%%%%%%%%%%%%%%%%%%%%%%%%%%%%%%%%%%%%%%%%%%%%%%%

Moments summarize important features of a probability distribution.

\begin{definitionbox}{Raw Moment}
The \(k\)-th raw moment of \(X\) is

\[
\boxed{
\mu_k'
=
\mathbb{E}[X^k],
}
\]

provided the expectation exists.
\end{definitionbox}

The first raw moment is

\[
\boxed{
\mu_1'
=
\mathbb{E}[X].
}
\]

%%%%%%%%%%%%%%%%%%%%%%%%%%%%%%%%%%%%%%%%%%%%%%%%%%%%%%%%%%%%
\subsection{Central Moments}
\label{subsec:central-moments}
%%%%%%%%%%%%%%%%%%%%%%%%%%%%%%%%%%%%%%%%%%%%%%%%%%%%%%%%%%%%

\begin{definitionbox}{Central Moment}
The \(k\)-th central moment of \(X\) is

\[
\boxed{
\mu_k
=
\mathbb{E}
\left[
(X-\mu)^k
\right],
}
\]

where

\[
\mu=\mathbb{E}[X].
\]
\end{definitionbox}

The first central moment is

\[
\boxed{
\mu_1=0.
}
\]

The second central moment is the variance.

%%%%%%%%%%%%%%%%%%%%%%%%%%%%%%%%%%%%%%%%%%%%%%%%%%%%%%%%%%%%
\subsection{Variance}
\label{subsec:variance}
%%%%%%%%%%%%%%%%%%%%%%%%%%%%%%%%%%%%%%%%%%%%%%%%%%%%%%%%%%%%

\begin{definitionbox}{Variance}
The variance of \(X\) is

\[
\boxed{
\operatorname{Var}(X)
=
\mathbb{E}
\left[
(X-\mathbb{E}[X])^2
\right].
}
\]
\end{definitionbox}

Variance measures average squared deviation from the mean.

It is always nonnegative.

%%%%%%%%%%%%%%%%%%%%%%%%%%%%%%%%%%%%%%%%%%%%%%%%%%%%%%%%%%%%
\subsection{Variance Is Nonnegative}
\label{subsec:variance-nonnegative}
%%%%%%%%%%%%%%%%%%%%%%%%%%%%%%%%%%%%%%%%%%%%%%%%%%%%%%%%%%%%

Since

\[
(X-\mathbb{E}[X])^2
\geq0,
\]

its expectation also satisfies

\[
\boxed{
\operatorname{Var}(X)\geq0.
}
\]

Moreover,

\[
\operatorname{Var}(X)=0
\]

if and only if \(X\) is constant almost surely.

%%%%%%%%%%%%%%%%%%%%%%%%%%%%%%%%%%%%%%%%%%%%%%%%%%%%%%%%%%%%
\subsection{Computational Formula for Variance}
\label{subsec:variance-computational-formula}
%%%%%%%%%%%%%%%%%%%%%%%%%%%%%%%%%%%%%%%%%%%%%%%%%%%%%%%%%%%%

Expand:

\[
(X-\mu)^2
=
X^2-2\mu X+\mu^2.
\]

Taking expectations,

\[
\operatorname{Var}(X)
=
\mathbb{E}[X^2]
-
2\mu\mathbb{E}[X]
+
\mu^2.
\]

Since

\[
\mu=\mathbb{E}[X],
\]

we obtain

\[
\boxed{
\operatorname{Var}(X)
=
\mathbb{E}[X^2]
-
(\mathbb{E}[X])^2.
}
\]

%%%%%%%%%%%%%%%%%%%%%%%%%%%%%%%%%%%%%%%%%%%%%%%%%%%%%%%%%%%%
\subsection{Worked Example: Variance of a Fair Die}
\label{subsec:worked-die-variance}
%%%%%%%%%%%%%%%%%%%%%%%%%%%%%%%%%%%%%%%%%%%%%%%%%%%%%%%%%%%%

\begin{workedexamplebox}{Variance of a Fair Die}

For a fair die,

\[
\mathbb{E}[X]
=
\frac72
\]

and

\[
\mathbb{E}[X^2]
=
\frac{91}{6}.
\]

Therefore,

\[
\operatorname{Var}(X)
=
\frac{91}{6}
-
\left(
\frac72
\right)^2.
\]

Thus,

\[
=
\frac{182}{12}
-
\frac{147}{12}.
\]

\begin{answerbox}
\[
\boxed{
\operatorname{Var}(X)
=
\frac{35}{12}.
}
\]
\end{answerbox}

\end{workedexamplebox}

%%%%%%%%%%%%%%%%%%%%%%%%%%%%%%%%%%%%%%%%%%%%%%%%%%%%%%%%%%%%
\subsection{Standard Deviation}
\label{subsec:standard-deviation}
%%%%%%%%%%%%%%%%%%%%%%%%%%%%%%%%%%%%%%%%%%%%%%%%%%%%%%%%%%%%

\begin{definitionbox}{Standard Deviation}
The standard deviation of \(X\) is

\[
\boxed{
\sigma_X
=
\sqrt{
\operatorname{Var}(X)
}.
}
\]
\end{definitionbox}

Variance has squared units.

Standard deviation has the same units as \(X\).

This often makes standard deviation easier to interpret.

%%%%%%%%%%%%%%%%%%%%%%%%%%%%%%%%%%%%%%%%%%%%%%%%%%%%%%%%%%%%
\subsection{Effect of Translation on Variance}
\label{subsec:translation-variance}
%%%%%%%%%%%%%%%%%%%%%%%%%%%%%%%%%%%%%%%%%%%%%%%%%%%%%%%%%%%%

Adding a constant does not change spread.

Thus,

\[
\boxed{
\operatorname{Var}(X+b)
=
\operatorname{Var}(X).
}
\]

The entire distribution shifts, but distances from its mean remain the
same.

%%%%%%%%%%%%%%%%%%%%%%%%%%%%%%%%%%%%%%%%%%%%%%%%%%%%%%%%%%%%
\subsection{Effect of Scaling on Variance}
\label{subsec:scaling-variance}
%%%%%%%%%%%%%%%%%%%%%%%%%%%%%%%%%%%%%%%%%%%%%%%%%%%%%%%%%%%%

For a constant \(a\),

\[
\boxed{
\operatorname{Var}(aX)
=
a^2\operatorname{Var}(X).
}
\]

Therefore,

\[
\boxed{
\operatorname{Var}(aX+b)
=
a^2\operatorname{Var}(X).
}
\]

Correspondingly,

\[
\boxed{
\sigma_{aX+b}
=
|a|\sigma_X.
}
\]

%%%%%%%%%%%%%%%%%%%%%%%%%%%%%%%%%%%%%%%%%%%%%%%%%%%%%%%%%%%%
\subsection{Worked Example: Linear Transformation}
\label{subsec:worked-linear-variance}
%%%%%%%%%%%%%%%%%%%%%%%%%%%%%%%%%%%%%%%%%%%%%%%%%%%%%%%%%%%%

\begin{workedexamplebox}{Transforming Mean and Variance}

Suppose

\[
\mathbb{E}[X]=8
\]

and

\[
\operatorname{Var}(X)=9.
\]

Let

\[
Y=4X-3.
\]

Then

\[
\mathbb{E}[Y]
=
4(8)-3
=
29.
\]

Also,

\[
\operatorname{Var}(Y)
=
4^2(9).
\]

\begin{answerbox}
\[
\boxed{
\mathbb{E}[Y]=29,
\qquad
\operatorname{Var}(Y)=144.
}
\]
\end{answerbox}

\end{workedexamplebox}

%%%%%%%%%%%%%%%%%%%%%%%%%%%%%%%%%%%%%%%%%%%%%%%%%%%%%%%%%%%%
\subsection{Mean and Variance of Bernoulli}
\label{subsec:bernoulli-mean-variance}
%%%%%%%%%%%%%%%%%%%%%%%%%%%%%%%%%%%%%%%%%%%%%%%%%%%%%%%%%%%%

If

\[
X\sim\operatorname{Bernoulli}(p),
\]

then

\[
\mathbb{E}[X]
=
0(1-p)+1(p).
\]

Thus,

\[
\boxed{
\mathbb{E}[X]=p.
}
\]

Since

\[
X^2=X
\]

for \(X\in\{0,1\}\),

\[
\mathbb{E}[X^2]=p.
\]

Therefore,

\[
\boxed{
\operatorname{Var}(X)
=
p(1-p).
}
\]

%%%%%%%%%%%%%%%%%%%%%%%%%%%%%%%%%%%%%%%%%%%%%%%%%%%%%%%%%%%%
\subsection{Mean and Variance of Binomial}
\label{subsec:binomial-mean-variance}
%%%%%%%%%%%%%%%%%%%%%%%%%%%%%%%%%%%%%%%%%%%%%%%%%%%%%%%%%%%%

If

\[
X\sim\operatorname{Bin}(n,p),
\]

write

\[
X
=
X_1+\cdots+X_n
\]

where the \(X_i\) are independent Bernoulli\((p)\) variables.

Then

\[
\mathbb{E}[X]
=
\sum_{i=1}^{n}
\mathbb{E}[X_i].
\]

Hence,

\[
\boxed{
\mathbb{E}[X]=np.
}
\]

Using independence,

\[
\boxed{
\operatorname{Var}(X)
=
np(1-p).
}
\]

%%%%%%%%%%%%%%%%%%%%%%%%%%%%%%%%%%%%%%%%%%%%%%%%%%%%%%%%%%%%
\subsection{Mean and Variance of Geometric}
\label{subsec:geometric-mean-variance}
%%%%%%%%%%%%%%%%%%%%%%%%%%%%%%%%%%%%%%%%%%%%%%%%%%%%%%%%%%%%

Under the convention that

\[
X
=
\text{number of trials until first success},
\]

if

\[
X\sim\operatorname{Geom}(p),
\]

then

\[
\boxed{
\mathbb{E}[X]
=
\frac1p.
}
\]

The variance is

\[
\boxed{
\operatorname{Var}(X)
=
\frac{1-p}{p^2}.
}
\]

%%%%%%%%%%%%%%%%%%%%%%%%%%%%%%%%%%%%%%%%%%%%%%%%%%%%%%%%%%%%
\subsection{Mean and Variance of Negative Binomial}
\label{subsec:negative-binomial-mean-variance}
%%%%%%%%%%%%%%%%%%%%%%%%%%%%%%%%%%%%%%%%%%%%%%%%%%%%%%%%%%%%

If

\[
X
\sim
\operatorname{NegBin}(r,p)
\]

counts the number of trials needed to obtain the \(r\)-th success, then

\[
\boxed{
\mathbb{E}[X]
=
\frac{r}{p}.
}
\]

Also,

\[
\boxed{
\operatorname{Var}(X)
=
\frac{r(1-p)}{p^2}.
}
\]

This follows by interpreting \(X\) as a sum of \(r\) independent
geometric waiting times.

%%%%%%%%%%%%%%%%%%%%%%%%%%%%%%%%%%%%%%%%%%%%%%%%%%%%%%%%%%%%
\subsection{Mean and Variance of Hypergeometric}
\label{subsec:hypergeometric-mean-variance}
%%%%%%%%%%%%%%%%%%%%%%%%%%%%%%%%%%%%%%%%%%%%%%%%%%%%%%%%%%%%

If

\[
X
\sim
\operatorname{Hypergeom}(N,K,n),
\]

then

\[
\boxed{
\mathbb{E}[X]
=
n\frac KN.
}
\]

The variance is

\[
\boxed{
\operatorname{Var}(X)
=
n
\frac KN
\left(
1-\frac KN
\right)
\frac{N-n}{N-1}.
}
\]

The factor

\[
\boxed{
\frac{N-n}{N-1}
}
\]

is called the finite population correction.

%%%%%%%%%%%%%%%%%%%%%%%%%%%%%%%%%%%%%%%%%%%%%%%%%%%%%%%%%%%%
\subsection{Mean and Variance of Poisson}
\label{subsec:poisson-mean-variance}
%%%%%%%%%%%%%%%%%%%%%%%%%%%%%%%%%%%%%%%%%%%%%%%%%%%%%%%%%%%%

If

\[
X\sim\operatorname{Poisson}(\lambda),
\]

then

\[
\boxed{
\mathbb{E}[X]
=
\lambda.
}
\]

Also,

\[
\boxed{
\operatorname{Var}(X)
=
\lambda.
}
\]

Thus the Poisson distribution has equal mean and variance.

%%%%%%%%%%%%%%%%%%%%%%%%%%%%%%%%%%%%%%%%%%%%%%%%%%%%%%%%%%%%
\subsection{Worked Example: Poisson Standard Deviation}
\label{subsec:worked-poisson-standard-deviation}
%%%%%%%%%%%%%%%%%%%%%%%%%%%%%%%%%%%%%%%%%%%%%%%%%%%%%%%%%%%%

\begin{workedexamplebox}{Poisson Spread}

Suppose

\[
X\sim\operatorname{Poisson}(16).
\]

Then

\[
\mathbb{E}[X]=16
\]

and

\[
\operatorname{Var}(X)=16.
\]

Therefore,

\begin{answerbox}
\[
\boxed{
\sigma_X=4.
}
\]
\end{answerbox}

\end{workedexamplebox}

%%%%%%%%%%%%%%%%%%%%%%%%%%%%%%%%%%%%%%%%%%%%%%%%%%%%%%%%%%%%
\subsection{Mean and Variance of Discrete Uniform}
\label{subsec:discrete-uniform-mean-variance}
%%%%%%%%%%%%%%%%%%%%%%%%%%%%%%%%%%%%%%%%%%%%%%%%%%%%%%%%%%%%

If

\[
X
\]

is uniformly distributed on

\[
\{1,2,\ldots,n\},
\]

then

\[
\boxed{
\mathbb{E}[X]
=
\frac{n+1}{2}.
}
\]

Its variance is

\[
\boxed{
\operatorname{Var}(X)
=
\frac{n^2-1}{12}.
}
\]

%%%%%%%%%%%%%%%%%%%%%%%%%%%%%%%%%%%%%%%%%%%%%%%%%%%%%%%%%%%%
\subsection{Mean and Variance of Continuous Uniform}
\label{subsec:continuous-uniform-mean-variance}
%%%%%%%%%%%%%%%%%%%%%%%%%%%%%%%%%%%%%%%%%%%%%%%%%%%%%%%%%%%%

If

\[
X\sim\operatorname{Uniform}(a,b),
\]

then

\[
\boxed{
\mathbb{E}[X]
=
\frac{a+b}{2}.
}
\]

The variance is

\[
\boxed{
\operatorname{Var}(X)
=
\frac{(b-a)^2}{12}.
}
\]

%%%%%%%%%%%%%%%%%%%%%%%%%%%%%%%%%%%%%%%%%%%%%%%%%%%%%%%%%%%%
\subsection{Mean and Variance of Exponential}
\label{subsec:exponential-mean-variance}
%%%%%%%%%%%%%%%%%%%%%%%%%%%%%%%%%%%%%%%%%%%%%%%%%%%%%%%%%%%%

If

\[
X\sim\operatorname{Exp}(\lambda),
\]

then

\[
\boxed{
\mathbb{E}[X]
=
\frac1\lambda.
}
\]

The variance is

\[
\boxed{
\operatorname{Var}(X)
=
\frac1{\lambda^2}.
}
\]

Thus,

\[
\boxed{
\sigma_X
=
\frac1\lambda.
}
\]

%%%%%%%%%%%%%%%%%%%%%%%%%%%%%%%%%%%%%%%%%%%%%%%%%%%%%%%%%%%%
\subsection{Mean and Variance of Gamma}
\label{subsec:gamma-mean-variance}
%%%%%%%%%%%%%%%%%%%%%%%%%%%%%%%%%%%%%%%%%%%%%%%%%%%%%%%%%%%%

Using the shape-rate parameterization,

\[
X
\sim
\operatorname{Gamma}(\alpha,\lambda),
\]

we have

\[
\boxed{
\mathbb{E}[X]
=
\frac{\alpha}{\lambda}.
}
\]

The variance is

\[
\boxed{
\operatorname{Var}(X)
=
\frac{\alpha}{\lambda^2}.
}
\]

%%%%%%%%%%%%%%%%%%%%%%%%%%%%%%%%%%%%%%%%%%%%%%%%%%%%%%%%%%%%
\subsection{Mean and Variance of Normal}
\label{subsec:normal-mean-variance}
%%%%%%%%%%%%%%%%%%%%%%%%%%%%%%%%%%%%%%%%%%%%%%%%%%%%%%%%%%%%

If

\[
X\sim N(\mu,\sigma^2),
\]

then by definition of the parameters,

\[
\boxed{
\mathbb{E}[X]
=
\mu
}
\]

and

\[
\boxed{
\operatorname{Var}(X)
=
\sigma^2.
}
\]

Thus the normal notation directly records its mean and variance.

%%%%%%%%%%%%%%%%%%%%%%%%%%%%%%%%%%%%%%%%%%%%%%%%%%%%%%%%%%%%
\subsection{Mean and Variance of Beta}
\label{subsec:beta-mean-variance}
%%%%%%%%%%%%%%%%%%%%%%%%%%%%%%%%%%%%%%%%%%%%%%%%%%%%%%%%%%%%

If

\[
X\sim\operatorname{Beta}(\alpha,\beta),
\]

then

\[
\boxed{
\mathbb{E}[X]
=
\frac{\alpha}{\alpha+\beta}.
}
\]

The variance is

\[
\boxed{
\operatorname{Var}(X)
=
\frac{
\alpha\beta
}{
(\alpha+\beta)^2
(\alpha+\beta+1)
}.
}
\]

%%%%%%%%%%%%%%%%%%%%%%%%%%%%%%%%%%%%%%%%%%%%%%%%%%%%%%%%%%%%
\subsection{Mean and Variance of Chi-Square}
\label{subsec:chi-square-mean-variance}
%%%%%%%%%%%%%%%%%%%%%%%%%%%%%%%%%%%%%%%%%%%%%%%%%%%%%%%%%%%%

If

\[
X\sim\chi^2_\nu,
\]

then

\[
\boxed{
\mathbb{E}[X]=\nu
}
\]

and

\[
\boxed{
\operatorname{Var}(X)=2\nu.
}
\]

These formulas follow from the gamma representation

\[
\chi^2_\nu
=
\operatorname{Gamma}
\left(
\frac{\nu}{2},
\frac12
\right).
\]

%%%%%%%%%%%%%%%%%%%%%%%%%%%%%%%%%%%%%%%%%%%%%%%%%%%%%%%%%%%%
\subsection{Mean of a Sum}
\label{subsec:mean-of-sum}
%%%%%%%%%%%%%%%%%%%%%%%%%%%%%%%%%%%%%%%%%%%%%%%%%%%%%%%%%%%%

For any finite collection of random variables with finite expectations,

\[
\boxed{
\mathbb{E}
\left[
\sum_{i=1}^{n}X_i
\right]
=
\sum_{i=1}^{n}
\mathbb{E}[X_i].
}
\]

This requires no independence.

Variance behaves differently.

%%%%%%%%%%%%%%%%%%%%%%%%%%%%%%%%%%%%%%%%%%%%%%%%%%%%%%%%%%%%
\subsection{Covariance}
\label{subsec:covariance}
%%%%%%%%%%%%%%%%%%%%%%%%%%%%%%%%%%%%%%%%%%%%%%%%%%%%%%%%%%%%

\begin{definitionbox}{Covariance}
For random variables \(X\) and \(Y\) with appropriate finite moments,

\[
\boxed{
\operatorname{Cov}(X,Y)
=
\mathbb{E}
\left[
(X-\mathbb{E}[X])
(Y-\mathbb{E}[Y])
\right].
}
\]
\end{definitionbox}

Covariance describes linear co-movement.

%%%%%%%%%%%%%%%%%%%%%%%%%%%%%%%%%%%%%%%%%%%%%%%%%%%%%%%%%%%%
\subsection{Interpretation of Covariance}
\label{subsec:covariance-interpretation}
%%%%%%%%%%%%%%%%%%%%%%%%%%%%%%%%%%%%%%%%%%%%%%%%%%%%%%%%%%%%

If \(X\) tends to be above its mean when \(Y\) is also above its mean,
covariance tends to be positive.

If \(X\) tends to be above its mean when \(Y\) is below its mean,
covariance tends to be negative.

If there is little linear co-movement, covariance may be near zero.

Thus:

\[
\boxed{
\operatorname{Cov}(X,Y)>0
}
\]

suggests positive linear association,

while

\[
\boxed{
\operatorname{Cov}(X,Y)<0
}
\]

suggests negative linear association.

%%%%%%%%%%%%%%%%%%%%%%%%%%%%%%%%%%%%%%%%%%%%%%%%%%%%%%%%%%%%
\subsection{Computational Formula for Covariance}
\label{subsec:covariance-computational}
%%%%%%%%%%%%%%%%%%%%%%%%%%%%%%%%%%%%%%%%%%%%%%%%%%%%%%%%%%%%

Expanding the definition gives

\[
\boxed{
\operatorname{Cov}(X,Y)
=
\mathbb{E}[XY]
-
\mathbb{E}[X]\mathbb{E}[Y].
}
\]

This is usually the most convenient computational formula.

%%%%%%%%%%%%%%%%%%%%%%%%%%%%%%%%%%%%%%%%%%%%%%%%%%%%%%%%%%%%
\subsection{Variance as Self-Covariance}
\label{subsec:variance-self-covariance}
%%%%%%%%%%%%%%%%%%%%%%%%%%%%%%%%%%%%%%%%%%%%%%%%%%%%%%%%%%%%

Setting

\[
Y=X,
\]

we obtain

\[
\boxed{
\operatorname{Cov}(X,X)
=
\operatorname{Var}(X).
}
\]

Thus variance is a special case of covariance.

%%%%%%%%%%%%%%%%%%%%%%%%%%%%%%%%%%%%%%%%%%%%%%%%%%%%%%%%%%%%
\subsection{Symmetry of Covariance}
\label{subsec:covariance-symmetry}
%%%%%%%%%%%%%%%%%%%%%%%%%%%%%%%%%%%%%%%%%%%%%%%%%%%%%%%%%%%%

Because multiplication is commutative,

\[
(X-\mu_X)(Y-\mu_Y)
=
(Y-\mu_Y)(X-\mu_X).
\]

Therefore,

\[
\boxed{
\operatorname{Cov}(X,Y)
=
\operatorname{Cov}(Y,X).
}
\]

%%%%%%%%%%%%%%%%%%%%%%%%%%%%%%%%%%%%%%%%%%%%%%%%%%%%%%%%%%%%
\subsection{Covariance with Constants}
\label{subsec:covariance-constants}
%%%%%%%%%%%%%%%%%%%%%%%%%%%%%%%%%%%%%%%%%%%%%%%%%%%%%%%%%%%%

For any constant \(c\),

\[
\boxed{
\operatorname{Cov}(X,c)=0.
}
\]

A constant does not vary, so it cannot co-vary with \(X\).

%%%%%%%%%%%%%%%%%%%%%%%%%%%%%%%%%%%%%%%%%%%%%%%%%%%%%%%%%%%%
\subsection{Scaling Covariance}
\label{subsec:scaling-covariance}
%%%%%%%%%%%%%%%%%%%%%%%%%%%%%%%%%%%%%%%%%%%%%%%%%%%%%%%%%%%%

For constants \(a,b,c,d\),

\[
\boxed{
\operatorname{Cov}(aX+b,cY+d)
=
ac\operatorname{Cov}(X,Y).
}
\]

Translations do not affect covariance.

Scalings multiply covariance.

%%%%%%%%%%%%%%%%%%%%%%%%%%%%%%%%%%%%%%%%%%%%%%%%%%%%%%%%%%%%
\subsection{Covariance of Sums}
\label{subsec:covariance-sums}
%%%%%%%%%%%%%%%%%%%%%%%%%%%%%%%%%%%%%%%%%%%%%%%%%%%%%%%%%%%%

Covariance is bilinear.

Thus,

\[
\boxed{
\operatorname{Cov}(X+Y,Z)
=
\operatorname{Cov}(X,Z)
+
\operatorname{Cov}(Y,Z).
}
\]

More generally,

\[
\boxed{
\operatorname{Cov}
\left(
\sum_i a_iX_i,
\sum_j b_jY_j
\right)
=
\sum_i\sum_j
a_ib_j
\operatorname{Cov}(X_i,Y_j).
}
\]

%%%%%%%%%%%%%%%%%%%%%%%%%%%%%%%%%%%%%%%%%%%%%%%%%%%%%%%%%%%%
\subsection{Variance of a Sum}
\label{subsec:variance-sum}
%%%%%%%%%%%%%%%%%%%%%%%%%%%%%%%%%%%%%%%%%%%%%%%%%%%%%%%%%%%%

Using covariance,

\[
\operatorname{Var}(X+Y)
=
\operatorname{Var}(X)
+
\operatorname{Var}(Y)
+
2\operatorname{Cov}(X,Y).
\]

Thus,

\[
\boxed{
\operatorname{Var}(X+Y)
=
\operatorname{Var}(X)
+
\operatorname{Var}(Y)
+
2\operatorname{Cov}(X,Y).
}
\]

%%%%%%%%%%%%%%%%%%%%%%%%%%%%%%%%%%%%%%%%%%%%%%%%%%%%%%%%%%%%
\subsection{Variance of a Difference}
\label{subsec:variance-difference}
%%%%%%%%%%%%%%%%%%%%%%%%%%%%%%%%%%%%%%%%%%%%%%%%%%%%%%%%%%%%

Similarly,

\[
\boxed{
\operatorname{Var}(X-Y)
=
\operatorname{Var}(X)
+
\operatorname{Var}(Y)
-
2\operatorname{Cov}(X,Y).
}
\]

The sign changes because \(Y\) is multiplied by \(-1\).

%%%%%%%%%%%%%%%%%%%%%%%%%%%%%%%%%%%%%%%%%%%%%%%%%%%%%%%%%%%%
\subsection{Variance of a General Linear Combination}
\label{subsec:variance-linear-combination}
%%%%%%%%%%%%%%%%%%%%%%%%%%%%%%%%%%%%%%%%%%%%%%%%%%%%%%%%%%%%

For constants \(a_1,\ldots,a_n\),

\[
\boxed{
\operatorname{Var}
\left(
\sum_{i=1}^{n}a_iX_i
\right)
=
\sum_{i=1}^{n}
a_i^2
\operatorname{Var}(X_i)
+
2
\sum_{i<j}
a_ia_j
\operatorname{Cov}(X_i,X_j).
}
\]

This is the general variance expansion.

%%%%%%%%%%%%%%%%%%%%%%%%%%%%%%%%%%%%%%%%%%%%%%%%%%%%%%%%%%%%
\subsection{Variance of Independent Sums}
\label{subsec:variance-independent-sums}
%%%%%%%%%%%%%%%%%%%%%%%%%%%%%%%%%%%%%%%%%%%%%%%%%%%%%%%%%%%%

If \(X\) and \(Y\) are independent and have finite second moments, then

\[
\operatorname{Cov}(X,Y)=0.
\]

Therefore,

\[
\boxed{
\operatorname{Var}(X+Y)
=
\operatorname{Var}(X)
+
\operatorname{Var}(Y).
}
\]

For independent \(X_1,\ldots,X_n\),

\[
\boxed{
\operatorname{Var}
\left(
\sum_{i=1}^{n}X_i
\right)
=
\sum_{i=1}^{n}
\operatorname{Var}(X_i).
}
\]

%%%%%%%%%%%%%%%%%%%%%%%%%%%%%%%%%%%%%%%%%%%%%%%%%%%%%%%%%%%%
\subsection{Independence Implies Zero Covariance}
\label{subsec:independence-zero-covariance}
%%%%%%%%%%%%%%%%%%%%%%%%%%%%%%%%%%%%%%%%%%%%%%%%%%%%%%%%%%%%

If \(X\) and \(Y\) are independent, then

\[
\mathbb{E}[XY]
=
\mathbb{E}[X]
\mathbb{E}[Y].
\]

Therefore,

\[
\boxed{
\operatorname{Cov}(X,Y)=0.
}
\]

Thus,

\[
\boxed{
\text{independence}
\Longrightarrow
\text{uncorrelatedness}.
}
\]

The converse does not generally hold.

%%%%%%%%%%%%%%%%%%%%%%%%%%%%%%%%%%%%%%%%%%%%%%%%%%%%%%%%%%%%
\subsection{Worked Example: Uncorrelated but Dependent}
\label{subsec:worked-uncorrelated-dependent}
%%%%%%%%%%%%%%%%%%%%%%%%%%%%%%%%%%%%%%%%%%%%%%%%%%%%%%%%%%%%

\begin{workedexamplebox}{A Nonlinear Dependence}

Let

\[
X\sim\operatorname{Uniform}(-1,1)
\]

and define

\[
Y=X^2.
\]

Then \(Y\) is completely determined by \(X\), so \(X\) and \(Y\) are
dependent.

By symmetry,

\[
\mathbb{E}[X]=0.
\]

Also,

\[
\mathbb{E}[XY]
=
\mathbb{E}[X^3]
=
0.
\]

Hence,

\[
\operatorname{Cov}(X,Y)
=
0.
\]

\begin{answerbox}
\[
\boxed{
X\text{ and }Y\text{ are dependent but uncorrelated}.
}
\]
\end{answerbox}

\end{workedexamplebox}

%%%%%%%%%%%%%%%%%%%%%%%%%%%%%%%%%%%%%%%%%%%%%%%%%%%%%%%%%%%%
\subsection{Correlation}
\label{subsec:correlation}
%%%%%%%%%%%%%%%%%%%%%%%%%%%%%%%%%%%%%%%%%%%%%%%%%%%%%%%%%%%%

\begin{definitionbox}{Correlation}
If

\[
\sigma_X>0
\]

and

\[
\sigma_Y>0,
\]

the correlation coefficient is

\[
\boxed{
\rho_{X,Y}
=
\frac{
\operatorname{Cov}(X,Y)
}{
\sigma_X\sigma_Y
}.
}
\]
\end{definitionbox}

The Greek letter

\[
\rho
\]

is called ``rho.''

%%%%%%%%%%%%%%%%%%%%%%%%%%%%%%%%%%%%%%%%%%%%%%%%%%%%%%%%%%%%
\subsection{Correlation Bounds}
\label{subsec:correlation-bounds-full}
%%%%%%%%%%%%%%%%%%%%%%%%%%%%%%%%%%%%%%%%%%%%%%%%%%%%%%%%%%%%

By the Cauchy--Schwarz inequality,

\[
\boxed{
-1
\leq
\rho_{X,Y}
\leq
1.
}
\]

Thus correlation is standardized to a dimensionless scale.

%%%%%%%%%%%%%%%%%%%%%%%%%%%%%%%%%%%%%%%%%%%%%%%%%%%%%%%%%%%%
\subsection{Meaning of Extreme Correlations}
\label{subsec:extreme-correlations}
%%%%%%%%%%%%%%%%%%%%%%%%%%%%%%%%%%%%%%%%%%%%%%%%%%%%%%%%%%%%

Under suitable nondegeneracy conditions,

\[
\rho_{X,Y}=1
\]

means there is a positive exact linear relationship

\[
Y=aX+b,
\qquad
a>0.
\]

Similarly,

\[
\rho_{X,Y}=-1
\]

corresponds to

\[
Y=aX+b,
\qquad
a<0.
\]

Thus perfect correlation represents exact linear dependence.

%%%%%%%%%%%%%%%%%%%%%%%%%%%%%%%%%%%%%%%%%%%%%%%%%%%%%%%%%%%%
\subsection{Correlation Is Invariant to Positive Rescaling}
\label{subsec:correlation-scale}
%%%%%%%%%%%%%%%%%%%%%%%%%%%%%%%%%%%%%%%%%%%%%%%%%%%%%%%%%%%%

Suppose

\[
U=aX+b
\]

and

\[
V=cY+d
\]

with

\[
a>0,
\qquad
c>0.
\]

Then

\[
\boxed{
\rho_{U,V}
=
\rho_{X,Y}.
}
\]

Changing measurement units therefore does not alter correlation.

If one scale factor is negative, the sign of correlation reverses.

%%%%%%%%%%%%%%%%%%%%%%%%%%%%%%%%%%%%%%%%%%%%%%%%%%%%%%%%%%%%
\subsection{Worked Example: Correlation from Covariance}
\label{subsec:worked-correlation}
%%%%%%%%%%%%%%%%%%%%%%%%%%%%%%%%%%%%%%%%%%%%%%%%%%%%%%%%%%%%

\begin{workedexamplebox}{Standardizing Covariance}

Suppose

\[
\operatorname{Cov}(X,Y)=6,
\]

\[
\sigma_X=2,
\]

and

\[
\sigma_Y=5.
\]

Then

\[
\rho_{X,Y}
=
\frac6{(2)(5)}.
\]

\begin{answerbox}
\[
\boxed{
\rho_{X,Y}=0.6.
}
\]
\end{answerbox}

\end{workedexamplebox}

%%%%%%%%%%%%%%%%%%%%%%%%%%%%%%%%%%%%%%%%%%%%%%%%%%%%%%%%%%%%
\subsection{Covariance Matrix}
\label{subsec:covariance-matrix}
%%%%%%%%%%%%%%%%%%%%%%%%%%%%%%%%%%%%%%%%%%%%%%%%%%%%%%%%%%%%

For a random vector

\[
\mathbf{X}
=
(X_1,\ldots,X_n)^T,
\]

define the covariance matrix

\[
\boxed{
\Sigma
=
\operatorname{Cov}(\mathbf{X}).
}
\]

Its entries are

\[
\boxed{
\Sigma_{ij}
=
\operatorname{Cov}(X_i,X_j).
}
\]

Thus,

\[
\Sigma
=
\begin{pmatrix}
\operatorname{Var}(X_1)
&
\operatorname{Cov}(X_1,X_2)
&
\cdots
\\
\operatorname{Cov}(X_2,X_1)
&
\operatorname{Var}(X_2)
&
\cdots
\\
\vdots
&
\vdots
&
\ddots
\end{pmatrix}.
\]

%%%%%%%%%%%%%%%%%%%%%%%%%%%%%%%%%%%%%%%%%%%%%%%%%%%%%%%%%%%%
\subsection{Properties of Covariance Matrices}
\label{subsec:covariance-matrix-properties}
%%%%%%%%%%%%%%%%%%%%%%%%%%%%%%%%%%%%%%%%%%%%%%%%%%%%%%%%%%%%

A covariance matrix is symmetric:

\[
\boxed{
\Sigma^T=\Sigma.
}
\]

It is also positive semidefinite.

For every vector \(a\),

\[
\boxed{
a^T\Sigma a\geq0.
}
\]

Indeed,

\[
a^T\Sigma a
=
\operatorname{Var}(a^T\mathbf{X}).
\]

%%%%%%%%%%%%%%%%%%%%%%%%%%%%%%%%%%%%%%%%%%%%%%%%%%%%%%%%%%%%
\subsection{Variance of a Linear Combination in Matrix Form}
\label{subsec:variance-matrix-form}
%%%%%%%%%%%%%%%%%%%%%%%%%%%%%%%%%%%%%%%%%%%%%%%%%%%%%%%%%%%%

Let

\[
Y=a^T\mathbf{X}.
\]

Then

\[
\boxed{
\operatorname{Var}(Y)
=
a^T\Sigma a.
}
\]

This compact formula is fundamental in multivariate statistics,
portfolio theory, and optimization.

%%%%%%%%%%%%%%%%%%%%%%%%%%%%%%%%%%%%%%%%%%%%%%%%%%%%%%%%%%%%
\subsection{Expectation of Products}
\label{subsec:expectation-products}
%%%%%%%%%%%%%%%%%%%%%%%%%%%%%%%%%%%%%%%%%%%%%%%%%%%%%%%%%%%%

In general,

\[
\boxed{
\mathbb{E}[XY]
\neq
\mathbb{E}[X]\mathbb{E}[Y].
}
\]

The equality holds when \(X\) and \(Y\) are independent under suitable
integrability conditions.

Thus,

\[
\boxed{
X\perp Y
\Longrightarrow
\mathbb{E}[g(X)h(Y)]
=
\mathbb{E}[g(X)]
\mathbb{E}[h(Y)]
}
\]

whenever the relevant expectations exist.

%%%%%%%%%%%%%%%%%%%%%%%%%%%%%%%%%%%%%%%%%%%%%%%%%%%%%%%%%%%%
\subsection{Expectation with Joint PMFs}
\label{subsec:joint-expectation-discrete}
%%%%%%%%%%%%%%%%%%%%%%%%%%%%%%%%%%%%%%%%%%%%%%%%%%%%%%%%%%%%

For discrete \(X,Y\),

\[
\boxed{
\mathbb{E}[g(X,Y)]
=
\sum_x\sum_y
g(x,y)
p_{X,Y}(x,y).
}
\]

In particular,

\[
\boxed{
\mathbb{E}[XY]
=
\sum_x\sum_y
xy
p_{X,Y}(x,y).
}
\]

%%%%%%%%%%%%%%%%%%%%%%%%%%%%%%%%%%%%%%%%%%%%%%%%%%%%%%%%%%%%
\subsection{Expectation with Joint Densities}
\label{subsec:joint-expectation-continuous}
%%%%%%%%%%%%%%%%%%%%%%%%%%%%%%%%%%%%%%%%%%%%%%%%%%%%%%%%%%%%

For continuous \(X,Y\) with joint density,

\[
\boxed{
\mathbb{E}[g(X,Y)]
=
\int_{-\infty}^{\infty}
\int_{-\infty}^{\infty}
g(x,y)
f_{X,Y}(x,y)
\,dy\,dx.
}
\]

In particular,

\[
\boxed{
\mathbb{E}[XY]
=
\iint
xyf_{X,Y}(x,y)
\,dx\,dy.
}
\]

%%%%%%%%%%%%%%%%%%%%%%%%%%%%%%%%%%%%%%%%%%%%%%%%%%%%%%%%%%%%
\subsection{Conditional Expectation: Discrete Case}
\label{subsec:conditional-expectation-discrete}
%%%%%%%%%%%%%%%%%%%%%%%%%%%%%%%%%%%%%%%%%%%%%%%%%%%%%%%%%%%%

Suppose \(X\) and \(Y\) are discrete.

For a value \(y\) with

\[
P(Y=y)>0,
\]

define

\[
\boxed{
\mathbb{E}[X\mid Y=y]
=
\sum_x
x
P(X=x\mid Y=y).
}
\]

This is the expected value of the conditional distribution of \(X\).

%%%%%%%%%%%%%%%%%%%%%%%%%%%%%%%%%%%%%%%%%%%%%%%%%%%%%%%%%%%%
\subsection{Conditional Expectation: Continuous Case}
\label{subsec:conditional-expectation-continuous}
%%%%%%%%%%%%%%%%%%%%%%%%%%%%%%%%%%%%%%%%%%%%%%%%%%%%%%%%%%%%

If a conditional density exists,

\[
\boxed{
\mathbb{E}[X\mid Y=y]
=
\int_{-\infty}^{\infty}
x
f_{X\mid Y}(x\mid y)
\,dx.
}
\]

As \(y\) varies, this produces a function of \(y\).

%%%%%%%%%%%%%%%%%%%%%%%%%%%%%%%%%%%%%%%%%%%%%%%%%%%%%%%%%%%%
\subsection{Conditional Expectation as a Random Variable}
\label{subsec:conditional-expectation-rv}
%%%%%%%%%%%%%%%%%%%%%%%%%%%%%%%%%%%%%%%%%%%%%%%%%%%%%%%%%%%%

The expression

\[
\boxed{
\mathbb{E}[X\mid Y]
}
\]

is itself a random variable.

If

\[
m(y)
=
\mathbb{E}[X\mid Y=y],
\]

then

\[
\boxed{
\mathbb{E}[X\mid Y]
=
m(Y).
}
\]

Thus conditional expectation is a function of the conditioning random
variable.

%%%%%%%%%%%%%%%%%%%%%%%%%%%%%%%%%%%%%%%%%%%%%%%%%%%%%%%%%%%%
\subsection{Law of Total Expectation}
\label{subsec:law-total-expectation}
%%%%%%%%%%%%%%%%%%%%%%%%%%%%%%%%%%%%%%%%%%%%%%%%%%%%%%%%%%%%

\begin{theorem}[Law of Total Expectation]
If the relevant expectations exist, then

\[
\boxed{
\mathbb{E}[X]
=
\mathbb{E}
\left[
\mathbb{E}[X\mid Y]
\right].
}
\]
\end{theorem}

This is also called the tower property in a more general setting.

%%%%%%%%%%%%%%%%%%%%%%%%%%%%%%%%%%%%%%%%%%%%%%%%%%%%%%%%%%%%
\subsection{Discrete Total Expectation Formula}
\label{subsec:discrete-total-expectation}
%%%%%%%%%%%%%%%%%%%%%%%%%%%%%%%%%%%%%%%%%%%%%%%%%%%%%%%%%%%%

If \(Y\) is discrete,

\[
\boxed{
\mathbb{E}[X]
=
\sum_y
\mathbb{E}[X\mid Y=y]
P(Y=y).
}
\]

Thus the unconditional mean is a weighted average of conditional means.

%%%%%%%%%%%%%%%%%%%%%%%%%%%%%%%%%%%%%%%%%%%%%%%%%%%%%%%%%%%%
\subsection{Worked Example: Total Expectation}
\label{subsec:worked-total-expectation}
%%%%%%%%%%%%%%%%%%%%%%%%%%%%%%%%%%%%%%%%%%%%%%%%%%%%%%%%%%%%

\begin{workedexamplebox}{Two Types of Customers}

Suppose a customer is type \(A\) with probability

\[
0.6
\]

and type \(B\) with probability

\[
0.4.
\]

Suppose

\[
\mathbb{E}[X\mid A]=5
\]

and

\[
\mathbb{E}[X\mid B]=9.
\]

Then

\[
\mathbb{E}[X]
=
5(0.6)+9(0.4).
\]

Therefore,

\begin{answerbox}
\[
\boxed{
\mathbb{E}[X]=6.6.
}
\]
\end{answerbox}

\end{workedexamplebox}

%%%%%%%%%%%%%%%%%%%%%%%%%%%%%%%%%%%%%%%%%%%%%%%%%%%%%%%%%%%%
\subsection{Conditional Expectation and Independence}
\label{subsec:conditional-expectation-independence}
%%%%%%%%%%%%%%%%%%%%%%%%%%%%%%%%%%%%%%%%%%%%%%%%%%%%%%%%%%%%

If \(X\) and \(Y\) are independent, then conditioning on \(Y\) does not
change the distribution of \(X\).

Therefore,

\[
\boxed{
\mathbb{E}[X\mid Y]
=
\mathbb{E}[X]
}
\]

almost surely, provided \(\mathbb{E}[X]\) exists.

%%%%%%%%%%%%%%%%%%%%%%%%%%%%%%%%%%%%%%%%%%%%%%%%%%%%%%%%%%%%
\subsection{Law of Total Variance}
\label{subsec:law-total-variance}
%%%%%%%%%%%%%%%%%%%%%%%%%%%%%%%%%%%%%%%%%%%%%%%%%%%%%%%%%%%%

\begin{theorem}[Law of Total Variance]
If \(X\) has finite variance, then

\[
\boxed{
\operatorname{Var}(X)
=
\mathbb{E}
\left[
\operatorname{Var}(X\mid Y)
\right]
+
\operatorname{Var}
\left(
\mathbb{E}[X\mid Y]
\right).
}
\]
\end{theorem}

This decomposes total variability into:

\[
\boxed{
\text{average within-group variability}
}
\]

plus

\[
\boxed{
\text{between-group variability}.
}
\]

%%%%%%%%%%%%%%%%%%%%%%%%%%%%%%%%%%%%%%%%%%%%%%%%%%%%%%%%%%%%
\subsection{Interpretation of Total Variance}
\label{subsec:total-variance-interpretation}
%%%%%%%%%%%%%%%%%%%%%%%%%%%%%%%%%%%%%%%%%%%%%%%%%%%%%%%%%%%%

The first term,

\[
\mathbb{E}[\operatorname{Var}(X\mid Y)],
\]

measures uncertainty remaining after \(Y\) is known.

The second term,

\[
\operatorname{Var}(\mathbb{E}[X\mid Y]),
\]

measures how much the conditional means differ across values of \(Y\).

Thus variability can be decomposed according to explanatory structure.

%%%%%%%%%%%%%%%%%%%%%%%%%%%%%%%%%%%%%%%%%%%%%%%%%%%%%%%%%%%%
\subsection{Worked Example: Variance Decomposition}
\label{subsec:worked-total-variance}
%%%%%%%%%%%%%%%%%%%%%%%%%%%%%%%%%%%%%%%%%%%%%%%%%%%%%%%%%%%%

\begin{workedexamplebox}{Within and Between Variability}

Suppose

\[
P(Y=0)=P(Y=1)=\frac12.
\]

Assume

\[
\mathbb{E}[X\mid Y=0]=2,
\qquad
\mathbb{E}[X\mid Y=1]=6,
\]

and

\[
\operatorname{Var}(X\mid Y=0)
=
\operatorname{Var}(X\mid Y=1)
=
3.
\]

Then

\[
\mathbb{E}[\operatorname{Var}(X\mid Y)]
=
3.
\]

The conditional mean variable takes values \(2\) and \(6\) equally
likely, so its mean is \(4\) and variance is

\[
\frac{
(2-4)^2+(6-4)^2
}{2}
=
4.
\]

Therefore,

\begin{answerbox}
\[
\boxed{
\operatorname{Var}(X)
=
3+4
=
7.
}
\]
\end{answerbox}

\end{workedexamplebox}

%%%%%%%%%%%%%%%%%%%%%%%%%%%%%%%%%%%%%%%%%%%%%%%%%%%%%%%%%%%%
\subsection{Higher Moments}
\label{subsec:higher-moments}
%%%%%%%%%%%%%%%%%%%%%%%%%%%%%%%%%%%%%%%%%%%%%%%%%%%%%%%%%%%%

Higher moments describe aspects of shape beyond location and spread.

The third central moment is

\[
\boxed{
\mu_3
=
\mathbb{E}
\left[
(X-\mu)^3
\right].
}
\]

The fourth central moment is

\[
\boxed{
\mu_4
=
\mathbb{E}
\left[
(X-\mu)^4
\right].
}
\]

These lead to skewness and kurtosis.

%%%%%%%%%%%%%%%%%%%%%%%%%%%%%%%%%%%%%%%%%%%%%%%%%%%%%%%%%%%%
\subsection{Skewness}
\label{subsec:skewness}
%%%%%%%%%%%%%%%%%%%%%%%%%%%%%%%%%%%%%%%%%%%%%%%%%%%%%%%%%%%%

\begin{definitionbox}{Skewness}
If the required moments exist and \(\sigma>0\), the standardized
skewness is

\[
\boxed{
\gamma_1
=
\frac{
\mathbb{E}[(X-\mu)^3]
}{
\sigma^3
}.
}
\]
\end{definitionbox}

The Greek letter

\[
\gamma
\]

is called ``gamma.''

%%%%%%%%%%%%%%%%%%%%%%%%%%%%%%%%%%%%%%%%%%%%%%%%%%%%%%%%%%%%
\subsection{Interpretation of Skewness}
\label{subsec:skewness-interpretation}
%%%%%%%%%%%%%%%%%%%%%%%%%%%%%%%%%%%%%%%%%%%%%%%%%%%%%%%%%%%%

Roughly,

\[
\gamma_1>0
\]

suggests right-skewness,

while

\[
\gamma_1<0
\]

suggests left-skewness.

For a distribution symmetric about its mean, when the third moment
exists,

\[
\boxed{
\gamma_1=0.
}
\]

Zero skewness does not necessarily imply symmetry.

%%%%%%%%%%%%%%%%%%%%%%%%%%%%%%%%%%%%%%%%%%%%%%%%%%%%%%%%%%%%
\subsection{Kurtosis}
\label{subsec:kurtosis}
%%%%%%%%%%%%%%%%%%%%%%%%%%%%%%%%%%%%%%%%%%%%%%%%%%%%%%%%%%%%

\begin{definitionbox}{Kurtosis}
If the fourth moment exists, kurtosis is

\[
\boxed{
\gamma_2
=
\frac{
\mathbb{E}[(X-\mu)^4]
}{
\sigma^4
}.
}
\]
\end{definitionbox}

For a normal distribution,

\[
\boxed{
\gamma_2=3.
}
\]

%%%%%%%%%%%%%%%%%%%%%%%%%%%%%%%%%%%%%%%%%%%%%%%%%%%%%%%%%%%%
\subsection{Excess Kurtosis}
\label{subsec:excess-kurtosis}
%%%%%%%%%%%%%%%%%%%%%%%%%%%%%%%%%%%%%%%%%%%%%%%%%%%%%%%%%%%%

Excess kurtosis is defined by

\[
\boxed{
\gamma_2-3.
}
\]

Thus the normal distribution has excess kurtosis

\[
\boxed{
0.
}
\]

Kurtosis is strongly influenced by tail behavior and extreme
observations.

%%%%%%%%%%%%%%%%%%%%%%%%%%%%%%%%%%%%%%%%%%%%%%%%%%%%%%%%%%%%
\subsection{Factorial Moments}
\label{subsec:factorial-moments}
%%%%%%%%%%%%%%%%%%%%%%%%%%%%%%%%%%%%%%%%%%%%%%%%%%%%%%%%%%%%

For an integer-valued random variable \(X\), the \(k\)-th falling
factorial is

\[
(X)_k
=
X(X-1)\cdots(X-k+1).
\]

The corresponding factorial moment is

\[
\boxed{
\mathbb{E}[(X)_k].
}
\]

Factorial moments are especially useful for count distributions.

%%%%%%%%%%%%%%%%%%%%%%%%%%%%%%%%%%%%%%%%%%%%%%%%%%%%%%%%%%%%
\subsection{Poisson Factorial Moments}
\label{subsec:poisson-factorial-moments}
%%%%%%%%%%%%%%%%%%%%%%%%%%%%%%%%%%%%%%%%%%%%%%%%%%%%%%%%%%%%

If

\[
X\sim\operatorname{Poisson}(\lambda),
\]

then

\[
\boxed{
\mathbb{E}[(X)_k]
=
\lambda^k.
}
\]

This simple formula is one reason the Poisson distribution appears
naturally in combinatorial counting problems.

%%%%%%%%%%%%%%%%%%%%%%%%%%%%%%%%%%%%%%%%%%%%%%%%%%%%%%%%%%%%
\subsection{Worked Example: Second Poisson Moment}
\label{subsec:worked-poisson-second-moment}
%%%%%%%%%%%%%%%%%%%%%%%%%%%%%%%%%%%%%%%%%%%%%%%%%%%%%%%%%%%%

\begin{workedexamplebox}{Using a Factorial Moment}

For

\[
X\sim\operatorname{Poisson}(\lambda),
\]

we know

\[
\mathbb{E}[X(X-1)]
=
\lambda^2.
\]

Since

\[
X^2
=
X(X-1)+X,
\]

we obtain

\[
\mathbb{E}[X^2]
=
\lambda^2+\lambda.
\]

Therefore,

\[
\operatorname{Var}(X)
=
\mathbb{E}[X^2]
-
(\mathbb{E}[X])^2.
\]

Thus,

\begin{answerbox}
\[
\boxed{
\operatorname{Var}(X)=\lambda.
}
\]
\end{answerbox}

\end{workedexamplebox}

%%%%%%%%%%%%%%%%%%%%%%%%%%%%%%%%%%%%%%%%%%%%%%%%%%%%%%%%%%%%
\subsection{Moment Existence}
\label{subsec:moment-existence}
%%%%%%%%%%%%%%%%%%%%%%%%%%%%%%%%%%%%%%%%%%%%%%%%%%%%%%%%%%%%

A distribution may have some finite moments but not others.

For example, a heavy-tailed distribution may satisfy

\[
\mathbb{E}[|X|]<\infty
\]

while

\[
\mathbb{E}[X^2]=\infty.
\]

Then the mean exists but the variance is infinite.

Thus one should not assume all moments exist.

%%%%%%%%%%%%%%%%%%%%%%%%%%%%%%%%%%%%%%%%%%%%%%%%%%%%%%%%%%%%
\subsection{The Cauchy Distribution and Moments}
\label{subsec:cauchy-moments}
%%%%%%%%%%%%%%%%%%%%%%%%%%%%%%%%%%%%%%%%%%%%%%%%%%%%%%%%%%%%

For a standard Cauchy variable,

\[
f_X(x)
=
\frac1{\pi(1+x^2)}.
\]

The expected value

\[
\mathbb{E}[X]
\]

does not exist as an absolutely convergent integral.

Its variance also does not exist.

Thus symmetry alone does not guarantee a finite mean.

\begin{warningbox}
A distribution can be perfectly symmetric and still fail to have a
finite expectation.
\end{warningbox}

%%%%%%%%%%%%%%%%%%%%%%%%%%%%%%%%%%%%%%%%%%%%%%%%%%%%%%%%%%%%
\subsection{Pareto Moment Existence}
\label{subsec:pareto-moment-existence}
%%%%%%%%%%%%%%%%%%%%%%%%%%%%%%%%%%%%%%%%%%%%%%%%%%%%%%%%%%%%

For a Pareto distribution with shape parameter \(\alpha\),

\[
\mathbb{E}[X^k]
\]

is finite only when

\[
\boxed{
k<\alpha.
}
\]

Therefore:

\[
\alpha\leq1
\]

implies the mean is infinite,

while

\[
\alpha\leq2
\]

implies the second moment and variance fail to be finite.

%%%%%%%%%%%%%%%%%%%%%%%%%%%%%%%%%%%%%%%%%%%%%%%%%%%%%%%%%%%%
\subsection{Markov's Inequality}
\label{subsec:markov-inequality}
%%%%%%%%%%%%%%%%%%%%%%%%%%%%%%%%%%%%%%%%%%%%%%%%%%%%%%%%%%%%

Expectations can provide probability bounds without knowing the full
distribution.

\begin{theorem}[Markov's Inequality]
If

\[
X\geq0
\]

and

\[
a>0,
\]

then

\[
\boxed{
P(X\geq a)
\leq
\frac{
\mathbb{E}[X]
}{a}.
}
\]
\end{theorem}

%%%%%%%%%%%%%%%%%%%%%%%%%%%%%%%%%%%%%%%%%%%%%%%%%%%%%%%%%%%%
\subsection{Proof of Markov's Inequality}
\label{subsec:proof-markov}
%%%%%%%%%%%%%%%%%%%%%%%%%%%%%%%%%%%%%%%%%%%%%%%%%%%%%%%%%%%%

Since \(X\geq0\),

\[
X
\geq
a\mathbf{1}_{\{X\geq a\}}.
\]

Taking expectations gives

\[
\mathbb{E}[X]
\geq
a
\mathbb{E}
\left[
\mathbf{1}_{\{X\geq a\}}
\right].
\]

Using

\[
\mathbb{E}
\left[
\mathbf{1}_{\{X\geq a\}}
\right]
=
P(X\geq a),
\]

we obtain

\[
\boxed{
P(X\geq a)
\leq
\frac{\mathbb{E}[X]}{a}.
}
\]

%%%%%%%%%%%%%%%%%%%%%%%%%%%%%%%%%%%%%%%%%%%%%%%%%%%%%%%%%%%%
\subsection{Worked Example: Markov Bound}
\label{subsec:worked-markov}
%%%%%%%%%%%%%%%%%%%%%%%%%%%%%%%%%%%%%%%%%%%%%%%%%%%%%%%%%%%%

\begin{workedexamplebox}{Bounding a Large Value}

Suppose

\[
X\geq0
\]

and

\[
\mathbb{E}[X]=10.
\]

Then

\[
P(X\geq50)
\leq
\frac{10}{50}.
\]

\begin{answerbox}
\[
\boxed{
P(X\geq50)\leq0.2.
}
\]
\end{answerbox}

\end{workedexamplebox}

%%%%%%%%%%%%%%%%%%%%%%%%%%%%%%%%%%%%%%%%%%%%%%%%%%%%%%%%%%%%
\subsection{Chebyshev's Inequality}
\label{subsec:chebyshev-inequality}
%%%%%%%%%%%%%%%%%%%%%%%%%%%%%%%%%%%%%%%%%%%%%%%%%%%%%%%%%%%%

\begin{theorem}[Chebyshev's Inequality]
If \(X\) has finite mean \(\mu\) and variance \(\sigma^2\), then for
every \(a>0\),

\[
\boxed{
P(|X-\mu|\geq a)
\leq
\frac{\sigma^2}{a^2}.
}
\]
\end{theorem}

Setting

\[
a=k\sigma
\]

gives

\[
\boxed{
P(|X-\mu|\geq k\sigma)
\leq
\frac1{k^2}.
}
\]

%%%%%%%%%%%%%%%%%%%%%%%%%%%%%%%%%%%%%%%%%%%%%%%%%%%%%%%%%%%%
\subsection{Proof of Chebyshev's Inequality}
\label{subsec:proof-chebyshev}
%%%%%%%%%%%%%%%%%%%%%%%%%%%%%%%%%%%%%%%%%%%%%%%%%%%%%%%%%%%%

Apply Markov's inequality to the nonnegative random variable

\[
(X-\mu)^2.
\]

Then

\[
P(|X-\mu|\geq a)
=
P((X-\mu)^2\geq a^2).
\]

Therefore,

\[
\leq
\frac{
\mathbb{E}[(X-\mu)^2]
}{
a^2
}.
\]

Hence,

\[
\boxed{
P(|X-\mu|\geq a)
\leq
\frac{\sigma^2}{a^2}.
}
\]

%%%%%%%%%%%%%%%%%%%%%%%%%%%%%%%%%%%%%%%%%%%%%%%%%%%%%%%%%%%%
\subsection{Worked Example: Chebyshev Bound}
\label{subsec:worked-chebyshev}
%%%%%%%%%%%%%%%%%%%%%%%%%%%%%%%%%%%%%%%%%%%%%%%%%%%%%%%%%%%%

\begin{workedexamplebox}{Within Three Standard Deviations}

For any random variable with finite variance,

\[
P(|X-\mu|\geq3\sigma)
\leq
\frac19.
\]

Therefore,

\[
P(|X-\mu|<3\sigma)
\geq
1-\frac19.
\]

\begin{answerbox}
\[
\boxed{
P(|X-\mu|<3\sigma)
\geq
\frac89.
}
\]
\end{answerbox}

\end{workedexamplebox}

%%%%%%%%%%%%%%%%%%%%%%%%%%%%%%%%%%%%%%%%%%%%%%%%%%%%%%%%%%%%
\subsection{Chebyshev Versus the Normal Empirical Rule}
\label{subsec:chebyshev-normal-comparison}
%%%%%%%%%%%%%%%%%%%%%%%%%%%%%%%%%%%%%%%%%%%%%%%%%%%%%%%%%%%%

Chebyshev's inequality applies to essentially any distribution with
finite variance.

It therefore gives relatively conservative bounds.

For \(k=2\),

\[
P(|X-\mu|<2\sigma)
\geq
\frac34.
\]

For a normal distribution, the actual probability is approximately

\[
0.95.
\]

Thus distribution-specific information can provide much sharper bounds.

%%%%%%%%%%%%%%%%%%%%%%%%%%%%%%%%%%%%%%%%%%%%%%%%%%%%%%%%%%%%
\subsection{One-Sided Chebyshev Preview}
\label{subsec:cantelli-preview}
%%%%%%%%%%%%%%%%%%%%%%%%%%%%%%%%%%%%%%%%%%%%%%%%%%%%%%%%%%%%

A useful one-sided refinement is Cantelli's inequality.

For \(a>0\),

\[
\boxed{
P(X-\mu\geq a)
\leq
\frac{
\sigma^2
}{
\sigma^2+a^2
}.
}
\]

This can be substantially sharper than applying ordinary Chebyshev to a
one-sided event.

%%%%%%%%%%%%%%%%%%%%%%%%%%%%%%%%%%%%%%%%%%%%%%%%%%%%%%%%%%%%
\subsection{Jensen's Inequality Preview}
\label{subsec:jensen-preview}
%%%%%%%%%%%%%%%%%%%%%%%%%%%%%%%%%%%%%%%%%%%%%%%%%%%%%%%%%%%%

Expectation also interacts with convexity.

If \(\varphi\) is convex and the expectations exist, then

\[
\boxed{
\varphi(\mathbb{E}[X])
\leq
\mathbb{E}[\varphi(X)].
}
\]

This is Jensen's inequality.

The symbol

\[
\varphi
\]

is called ``phi'' or ``varphi.''

Jensen's inequality is developed more fully in analysis and advanced
probability.

%%%%%%%%%%%%%%%%%%%%%%%%%%%%%%%%%%%%%%%%%%%%%%%%%%%%%%%%%%%%
\subsection{Variance from Jensen's Inequality}
\label{subsec:variance-jensen}
%%%%%%%%%%%%%%%%%%%%%%%%%%%%%%%%%%%%%%%%%%%%%%%%%%%%%%%%%%%%

Take

\[
\varphi(x)=x^2,
\]

which is convex.

Then

\[
(\mathbb{E}[X])^2
\leq
\mathbb{E}[X^2].
\]

Therefore,

\[
\boxed{
\mathbb{E}[X^2]
-
(\mathbb{E}[X])^2
\geq0.
}
\]

This gives another proof that variance is nonnegative.

%%%%%%%%%%%%%%%%%%%%%%%%%%%%%%%%%%%%%%%%%%%%%%%%%%%%%%%%%%%%
\subsection{Moment Inequalities}
\label{subsec:moment-inequalities-preview}
%%%%%%%%%%%%%%%%%%%%%%%%%%%%%%%%%%%%%%%%%%%%%%%%%%%%%%%%%%%%

Several major inequalities connect moments:

\[
\boxed{
\text{Cauchy--Schwarz},
}
\]

\[
\boxed{
\text{Jensen},
}
\]

\[
\boxed{
\text{Hölder},
}
\]

and

\[
\boxed{
\text{Minkowski}.
}
\]

These become fundamental tools in measure-theoretic probability and
functional analysis.

%%%%%%%%%%%%%%%%%%%%%%%%%%%%%%%%%%%%%%%%%%%%%%%%%%%%%%%%%%%%
\subsection{Cauchy--Schwarz for Random Variables}
\label{subsec:cauchy-schwarz-rv}
%%%%%%%%%%%%%%%%%%%%%%%%%%%%%%%%%%%%%%%%%%%%%%%%%%%%%%%%%%%%

If \(X\) and \(Y\) have finite second moments, then

\[
\boxed{
|\mathbb{E}[XY]|
\leq
\sqrt{
\mathbb{E}[X^2]
\mathbb{E}[Y^2]
}.
}
\]

Applying this to centered variables gives

\[
\boxed{
|\operatorname{Cov}(X,Y)|
\leq
\sigma_X\sigma_Y.
}
\]

This proves the correlation bound.

%%%%%%%%%%%%%%%%%%%%%%%%%%%%%%%%%%%%%%%%%%%%%%%%%%%%%%%%%%%%
\subsection{Standardization}
\label{subsec:standardization-expectation}
%%%%%%%%%%%%%%%%%%%%%%%%%%%%%%%%%%%%%%%%%%%%%%%%%%%%%%%%%%%%

For a random variable with finite mean \(\mu\) and positive standard
deviation \(\sigma\), define

\[
\boxed{
Z
=
\frac{X-\mu}{\sigma}.
}
\]

Then

\[
\boxed{
\mathbb{E}[Z]=0.
}
\]

Also,

\[
\boxed{
\operatorname{Var}(Z)=1.
}
\]

Thus standardization centers and rescales a random variable.

%%%%%%%%%%%%%%%%%%%%%%%%%%%%%%%%%%%%%%%%%%%%%%%%%%%%%%%%%%%%
\subsection{Worked Example: Standardizing an Arbitrary Variable}
\label{subsec:worked-standardization-general}
%%%%%%%%%%%%%%%%%%%%%%%%%%%%%%%%%%%%%%%%%%%%%%%%%%%%%%%%%%%%

\begin{workedexamplebox}{Mean Zero and Variance One}

Suppose

\[
\mathbb{E}[X]=12
\]

and

\[
\operatorname{Var}(X)=25.
\]

Then

\[
\sigma_X=5.
\]

Define

\[
Z
=
\frac{X-12}{5}.
\]

Therefore,

\begin{answerbox}
\[
\boxed{
\mathbb{E}[Z]=0,
\qquad
\operatorname{Var}(Z)=1.
}
\]
\end{answerbox}

\end{workedexamplebox}

%%%%%%%%%%%%%%%%%%%%%%%%%%%%%%%%%%%%%%%%%%%%%%%%%%%%%%%%%%%%
\subsection{Expectation and Probability Generating Functions}
\label{subsec:expectation-pgf-preview}
%%%%%%%%%%%%%%%%%%%%%%%%%%%%%%%%%%%%%%%%%%%%%%%%%%%%%%%%%%%%

For a nonnegative integer-valued random variable with PGF

\[
G_X(s)
=
\mathbb{E}[s^X],
\]

differentiation gives moment information.

Under appropriate conditions,

\[
\boxed{
G_X'(1)
=
\mathbb{E}[X].
}
\]

Also,

\[
\boxed{
G_X''(1)
=
\mathbb{E}[X(X-1)].
}
\]

Thus PGFs naturally encode factorial moments.

%%%%%%%%%%%%%%%%%%%%%%%%%%%%%%%%%%%%%%%%%%%%%%%%%%%%%%%%%%%%
\subsection{Moment Generating Functions Preview}
\label{subsec:mgf-preview-expectation}
%%%%%%%%%%%%%%%%%%%%%%%%%%%%%%%%%%%%%%%%%%%%%%%%%%%%%%%%%%%%

The moment generating function is

\[
\boxed{
M_X(t)
=
\mathbb{E}[e^{tX}].
}
\]

When it exists in a neighborhood of \(t=0\),

\[
\boxed{
M_X^{(k)}(0)
=
\mathbb{E}[X^k].
}
\]

Thus derivatives of the MGF generate raw moments.

This is developed systematically in Section~\ref{sec:probability-transforms}.

%%%%%%%%%%%%%%%%%%%%%%%%%%%%%%%%%%%%%%%%%%%%%%%%%%%%%%%%%%%%
\subsection{Cumulants Preview}
\label{subsec:cumulants-preview}
%%%%%%%%%%%%%%%%%%%%%%%%%%%%%%%%%%%%%%%%%%%%%%%%%%%%%%%%%%%%

The cumulant generating function is

\[
\boxed{
K_X(t)
=
\ln M_X(t).
}
\]

Its derivatives at zero are called cumulants.

The first cumulant is the mean:

\[
\boxed{
\kappa_1
=
\mathbb{E}[X].
}
\]

The second is the variance:

\[
\boxed{
\kappa_2
=
\operatorname{Var}(X).
}
\]

Cumulants are especially useful for sums of independent random
variables.

%%%%%%%%%%%%%%%%%%%%%%%%%%%%%%%%%%%%%%%%%%%%%%%%%%%%%%%%%%%%
\subsection{Expectation as Center of Mass}
\label{subsec:expectation-center-mass}
%%%%%%%%%%%%%%%%%%%%%%%%%%%%%%%%%%%%%%%%%%%%%%%%%%%%%%%%%%%%

A discrete probability distribution can be interpreted as point masses.

Place mass

\[
p_X(x)
\]

at each location \(x\).

Then

\[
\boxed{
\mathbb{E}[X]
}
\]

is the center of mass.

For continuous distributions, the density plays the role of distributed
mass.

This geometric interpretation explains many expectation formulas.

%%%%%%%%%%%%%%%%%%%%%%%%%%%%%%%%%%%%%%%%%%%%%%%%%%%%%%%%%%%%
\subsection{Variance as Moment of Inertia Analogy}
\label{subsec:variance-inertia}
%%%%%%%%%%%%%%%%%%%%%%%%%%%%%%%%%%%%%%%%%%%%%%%%%%%%%%%%%%%%

Variance measures squared distance from the center.

Thus

\[
\boxed{
\operatorname{Var}(X)
=
\mathbb{E}[(X-\mu)^2]
}
\]

is analogous to a one-dimensional moment of inertia around the center
\(\mu\).

Far-away probability receives much more weight because deviations are
squared.

%%%%%%%%%%%%%%%%%%%%%%%%%%%%%%%%%%%%%%%%%%%%%%%%%%%%%%%%%%%%
\subsection{Expectation in Gambling and Fair Games}
\label{subsec:expectation-fair-games}
%%%%%%%%%%%%%%%%%%%%%%%%%%%%%%%%%%%%%%%%%%%%%%%%%%%%%%%%%%%%

Suppose a game produces net payoff \(X\).

The game is often called fair if

\[
\boxed{
\mathbb{E}[X]=0.
}
\]

A positive expected payoff favors the player in the long-run average
sense.

A negative expected payoff favors the opposing side.

Expected value alone does not measure risk.

Variance and the full distribution are also important.

%%%%%%%%%%%%%%%%%%%%%%%%%%%%%%%%%%%%%%%%%%%%%%%%%%%%%%%%%%%%
\subsection{Worked Example: Expected Game Payoff}
\label{subsec:worked-game-payoff}
%%%%%%%%%%%%%%%%%%%%%%%%%%%%%%%%%%%%%%%%%%%%%%%%%%%%%%%%%%%%

\begin{workedexamplebox}{A Simple Game}

Suppose a player wins \$10 with probability

\[
0.2
\]

and loses \$3 with probability

\[
0.8.
\]

Let \(X\) be net payoff.

Then

\[
\mathbb{E}[X]
=
10(0.2)
+
(-3)(0.8).
\]

Therefore,

\[
=
2-2.4.
\]

\begin{answerbox}
\[
\boxed{
\mathbb{E}[X]
=
-0.4.
}
\]
\end{answerbox}

The expected net loss is \$0.40 per play.
\end{workedexamplebox}

%%%%%%%%%%%%%%%%%%%%%%%%%%%%%%%%%%%%%%%%%%%%%%%%%%%%%%%%%%%%
\subsection{Expectation and Long-Run Behavior}
\label{subsec:expectation-long-run}
%%%%%%%%%%%%%%%%%%%%%%%%%%%%%%%%%%%%%%%%%%%%%%%%%%%%%%%%%%%%

Expected value is connected to long-run averages through laws of large
numbers.

If

\[
X_1,X_2,\ldots
\]

are suitable repeated observations with mean \(\mu\), then the sample
average

\[
\frac{
X_1+\cdots+X_n
}{n}
\]

tends toward \(\mu\) in an appropriate probabilistic sense.

This is developed later in the chapter.

%%%%%%%%%%%%%%%%%%%%%%%%%%%%%%%%%%%%%%%%%%%%%%%%%%%%%%%%%%%%
\subsection{Common Errors}
\label{subsec:expectation-moments-common-errors}
%%%%%%%%%%%%%%%%%%%%%%%%%%%%%%%%%%%%%%%%%%%%%%%%%%%%%%%%%%%%

\subsubsection{Assuming Expected Value Must Be a Possible Outcome}

Expectation is a weighted average.

It need not belong to the support.

%%%%%%%%%%%%%%%%%%%%%%%%%%%%%%%%%%%%%%%%%%%%%%%%%%%%%%%%%%%%
\subsubsection{Forgetting Probability Weights}

For a discrete random variable,

\[
\mathbb{E}[X]
\neq
\sum_x x.
\]

The correct expression is

\[
\boxed{
\mathbb{E}[X]
=
\sum_xxp_X(x).
}
\]

%%%%%%%%%%%%%%%%%%%%%%%%%%%%%%%%%%%%%%%%%%%%%%%%%%%%%%%%%%%%
\subsubsection{Assuming Expectation Always Exists}

Heavy-tailed distributions may have undefined or infinite moments.

Always check convergence when necessary.

%%%%%%%%%%%%%%%%%%%%%%%%%%%%%%%%%%%%%%%%%%%%%%%%%%%%%%%%%%%%
\subsubsection{Assuming Variance Is \(\mathbb{E}[X^2]\)}

The correct formula is

\[
\boxed{
\operatorname{Var}(X)
=
\mathbb{E}[X^2]
-
(\mathbb{E}[X])^2.
}
\]

%%%%%%%%%%%%%%%%%%%%%%%%%%%%%%%%%%%%%%%%%%%%%%%%%%%%%%%%%%%%
\subsubsection{Forgetting to Square Scale Factors in Variance}

For

\[
Y=aX+b,
\]

\[
\boxed{
\operatorname{Var}(Y)
=
a^2\operatorname{Var}(X).
}
\]

%%%%%%%%%%%%%%%%%%%%%%%%%%%%%%%%%%%%%%%%%%%%%%%%%%%%%%%%%%%%
\subsubsection{Adding Variances Without Checking Covariance}

In general,

\[
\operatorname{Var}(X+Y)
\neq
\operatorname{Var}(X)
+
\operatorname{Var}(Y).
\]

The covariance term must be included unless it is zero.

%%%%%%%%%%%%%%%%%%%%%%%%%%%%%%%%%%%%%%%%%%%%%%%%%%%%%%%%%%%%
\subsubsection{Assuming Zero Covariance Means Independence}

Zero covariance only rules out linear association.

Nonlinear dependence may remain.

%%%%%%%%%%%%%%%%%%%%%%%%%%%%%%%%%%%%%%%%%%%%%%%%%%%%%%%%%%%%
\subsubsection{Assuming Independence Is Needed for Linearity of Expectation}

Linearity of expectation does not require independence.

%%%%%%%%%%%%%%%%%%%%%%%%%%%%%%%%%%%%%%%%%%%%%%%%%%%%%%%%%%%%
\subsubsection{Confusing Variance and Standard Deviation}

Variance has squared units.

Standard deviation has the same units as the random variable.

%%%%%%%%%%%%%%%%%%%%%%%%%%%%%%%%%%%%%%%%%%%%%%%%%%%%%%%%%%%%
\subsubsection{Confusing Conditional Mean with Marginal Mean}

The value

\[
\mathbb{E}[X\mid Y=y]
\]

depends on the conditioning value \(y\).

The unconditional mean averages these conditional means over the
distribution of \(Y\).

%%%%%%%%%%%%%%%%%%%%%%%%%%%%%%%%%%%%%%%%%%%%%%%%%%%%%%%%%%%%
\subsubsection{Interpreting Correlation as Causation}

Correlation describes statistical association.

It does not by itself establish a causal relationship.

%%%%%%%%%%%%%%%%%%%%%%%%%%%%%%%%%%%%%%%%%%%%%%%%%%%%%%%%%%%%
\subsubsection{Assuming Small Variance Means Small Mean}

Mean and variance describe different features.

A distribution can have a large mean and small variance, or vice versa.

%%%%%%%%%%%%%%%%%%%%%%%%%%%%%%%%%%%%%%%%%%%%%%%%%%%%%%%%%%%%
\subsection{Applications}
\label{subsec:expectation-moments-applications}
%%%%%%%%%%%%%%%%%%%%%%%%%%%%%%%%%%%%%%%%%%%%%%%%%%%%%%%%%%%%

\begin{applicationbox}

\textbf{Statistics.}
Means, variances, covariances, and higher moments form the foundation of
estimation, sampling theory, regression, and statistical inference.

\medskip

\textbf{Finance.}
Expected return, variance, covariance matrices, and correlation are
central to portfolio and risk models.

\medskip

\textbf{Engineering.}
Expected load, expected lifetime, uncertainty, and reliability are
summarized using moments.

\medskip

\textbf{Operations Research.}
Expected demand, waiting time, cost, and service requirements guide
decision models.

\medskip

\textbf{Combinatorics.}
Indicator variables and linearity of expectation provide powerful
methods for counting random structures.

\medskip

\textbf{Machine Learning.}
Means and covariance matrices summarize distributions and appear in
Gaussian models, optimization, and dimensionality reduction.

\medskip

\textbf{Physics.}
Expectation values, fluctuations, and covariance describe uncertain
physical quantities.

\medskip

\textbf{Biology.}
Means, variances, and correlations summarize biological measurements,
population variation, and experimental outcomes.

\end{applicationbox}

%%%%%%%%%%%%%%%%%%%%%%%%%%%%%%%%%%%%%%%%%%%%%%%%%%%%%%%%%%%%
\subsection{Exercises}
\label{subsec:expectation-moments-exercises}
%%%%%%%%%%%%%%%%%%%%%%%%%%%%%%%%%%%%%%%%%%%%%%%%%%%%%%%%%%%%

\begin{exercise}[1]
Define the expected value of a discrete random variable.
\end{exercise}

\begin{exercise}[1]
Define the expected value of a continuous random variable.
\end{exercise}

\begin{exercise}[1]
Define the \(k\)-th raw moment.
\end{exercise}

\begin{exercise}[1]
Define the \(k\)-th central moment.
\end{exercise}

\begin{exercise}[1]
Define variance.
\end{exercise}

\begin{exercise}[1]
State the computational formula for variance.
\end{exercise}

\begin{exercise}[1]
Define standard deviation.
\end{exercise}

\begin{exercise}[1]
Define covariance.
\end{exercise}

\begin{exercise}[1]
Define correlation.
\end{exercise}

\begin{exercise}[1]
State the law of total expectation.
\end{exercise}

\begin{exercise}[1]
State the law of total variance.
\end{exercise}

\begin{exercise}[2]
Let \(X\) have PMF

\[
P(X=0)=0.2,
\qquad
P(X=1)=0.5,
\qquad
P(X=2)=0.3.
\]

Find

\[
\mathbb{E}[X].
\]
\end{exercise}

\begin{exercise}[2]
For the same variable, find

\[
\mathbb{E}[X^2].
\]
\end{exercise}

\begin{exercise}[2]
Use the previous results to find

\[
\operatorname{Var}(X).
\]
\end{exercise}

\begin{exercise}[2]
Suppose

\[
\mathbb{E}[X]=7
\]

and

\[
\operatorname{Var}(X)=4.
\]

Find the mean and variance of

\[
Y=3X-5.
\]
\end{exercise}

\begin{exercise}[2]
Let

\[
X\sim\operatorname{Bernoulli}(0.4).
\]

Find its mean and variance.
\end{exercise}

\begin{exercise}[2]
Let

\[
X\sim\operatorname{Bin}(20,0.3).
\]

Find its mean, variance, and standard deviation.
\end{exercise}

\begin{exercise}[2]
Let

\[
X\sim\operatorname{Poisson}(12).
\]

Find its mean and variance.
\end{exercise}

\begin{exercise}[2]
Let

\[
X\sim\operatorname{Exp}(4).
\]

Find its mean and variance.
\end{exercise}

\begin{exercise}[2]
Let

\[
X\sim\operatorname{Uniform}(2,14).
\]

Find its mean and variance.
\end{exercise}

\begin{exercise}[2]
Let

\[
X\sim N(50,9).
\]

Find the mean and standard deviation.
\end{exercise}

\begin{exercise}[2]
Suppose

\[
\operatorname{Cov}(X,Y)=8,
\qquad
\sigma_X=4,
\qquad
\sigma_Y=5.
\]

Find the correlation.
\end{exercise}

\begin{exercise}[2]
Suppose

\[
\operatorname{Var}(X)=9,
\]

\[
\operatorname{Var}(Y)=16,
\]

and

\[
\operatorname{Cov}(X,Y)=3.
\]

Find

\[
\operatorname{Var}(X+Y).
\]
\end{exercise}

\begin{exercise}[2]
Using the same values, find

\[
\operatorname{Var}(X-Y).
\]
\end{exercise}

\begin{exercise}[2]
Suppose \(X\) and \(Y\) are independent with variances \(4\) and \(7\).

Find

\[
\operatorname{Var}(2X-3Y).
\]
\end{exercise}

\begin{exercise}[3]
Prove linearity of expectation for two discrete random variables.
\end{exercise}

\begin{exercise}[3]
Prove

\[
\mathbb{E}[aX+b]
=
a\mathbb{E}[X]+b.
\]
\end{exercise}

\begin{exercise}[3]
Prove

\[
\operatorname{Var}(X)
=
\mathbb{E}[X^2]
-
(\mathbb{E}[X])^2.
\]
\end{exercise}

\begin{exercise}[3]
Prove

\[
\operatorname{Var}(aX+b)
=
a^2\operatorname{Var}(X).
\]
\end{exercise}

\begin{exercise}[3]
Prove

\[
\operatorname{Cov}(X,Y)
=
\mathbb{E}[XY]
-
\mathbb{E}[X]\mathbb{E}[Y].
\]
\end{exercise}

\begin{exercise}[3]
Derive

\[
\operatorname{Var}(X+Y)
=
\operatorname{Var}(X)
+
\operatorname{Var}(Y)
+
2\operatorname{Cov}(X,Y).
\]
\end{exercise}

\begin{exercise}[3]
Show that independence implies zero covariance.
\end{exercise}

\begin{exercise}[3]
Construct an example of dependent random variables with zero covariance.
\end{exercise}

\begin{exercise}[3]
Prove that

\[
\mathbb{E}[\mathbf{1}_A]
=
P(A).
\]
\end{exercise}

\begin{exercise}[3]
Use indicator variables to find the expected number of heads in \(n\)
fair coin tosses.
\end{exercise}

\begin{exercise}[3]
Use indicator variables to find the expected number of fixed points in a
uniformly random permutation of \(n\) objects.
\end{exercise}

\begin{exercise}[3]
Use indicator variables to find the expected number of edges in the
random graph \(G(n,p)\).
\end{exercise}

\begin{exercise}[3]
Use indicator variables to find the expected number of triangles in
\(G(n,p)\).
\end{exercise}

\begin{exercise}[3]
Derive the mean of the binomial distribution using indicators.
\end{exercise}

\begin{exercise}[3]
Derive the variance of the binomial distribution using independent
Bernoulli variables.
\end{exercise}

\begin{exercise}[3]
Derive the mean of the exponential distribution directly from its
density.
\end{exercise}

\begin{exercise}[3]
Derive the variance of the exponential distribution.
\end{exercise}

\begin{exercise}[3]
Prove the tail-sum formula for a nonnegative integer-valued random
variable.
\end{exercise}

\begin{exercise}[3]
Prove the continuous tail formula

\[
\mathbb{E}[X]
=
\int_0^\infty P(X>x)\,dx
\]

for a nonnegative continuous random variable.
\end{exercise}

\begin{exercise}[3]
Prove Markov's inequality.
\end{exercise}

\begin{exercise}[3]
Derive Chebyshev's inequality from Markov's inequality.
\end{exercise}

\begin{exercise}[3]
Use Chebyshev's inequality to show that at least \(15/16\) of the
probability lies within \(4\) standard deviations of the mean.
\end{exercise}

\begin{exercise}[3]
Prove the law of total expectation for a discrete conditioning random
variable.
\end{exercise}

\begin{exercise}[3]
Derive the law of total variance.
\end{exercise}

\begin{exercise}[3]
Show that every covariance matrix is symmetric.
\end{exercise}

\begin{exercise}[3]
Show that every covariance matrix is positive semidefinite.
\end{exercise}

\begin{exercise}[3]
Prove

\[
\operatorname{Var}(a^T\mathbf{X})
=
a^T\Sigma a.
\]
\end{exercise}

\begin{exercise}[4]
Study Jensen's inequality and prove it for convex functions under
appropriate conditions.
\end{exercise}

\begin{exercise}[4]
Study Hölder's inequality and derive Cauchy--Schwarz as a special case.
\end{exercise}

\begin{exercise}[4]
Study Minkowski's inequality and its relationship with \(L^p\) norms.
\end{exercise}

\begin{exercise}[4]
Investigate Cantelli's one-sided variance inequality.
\end{exercise}

\begin{exercise}[4]
Study higher-order central moments of the normal distribution.
\end{exercise}

\begin{exercise}[4]
Derive skewness and kurtosis formulas for the gamma distribution.
\end{exercise}

\begin{exercise}[4]
Investigate moment existence for the Pareto distribution.
\end{exercise}

\begin{exercise}[4]
Explain carefully why the standard Cauchy distribution has no finite
expectation.
\end{exercise}

\begin{exercise}[4]
Study cumulants and show that cumulants add for independent random
variables.
\end{exercise}

\begin{exercise}[4]
Investigate the method of moments in statistical estimation.
\end{exercise}

\begin{exercise}[4]
Study factorial moments of Poisson and binomial variables.
\end{exercise}

\begin{exercise}[4]
Investigate conditional expectation as an orthogonal projection in
\(L^2\).
\end{exercise}

\begin{exercise}[4]
Study covariance operators for infinite-dimensional random objects.
\end{exercise}

%%%%%%%%%%%%%%%%%%%%%%%%%%%%%%%%%%%%%%%%%%%%%%%%%%%%%%%%%%%%
\subsection{Conceptual Summary}
\label{subsec:expectation-moments-summary}
%%%%%%%%%%%%%%%%%%%%%%%%%%%%%%%%%%%%%%%%%%%%%%%%%%%%%%%%%%%%

For a discrete random variable,

\[
\boxed{
\mathbb{E}[X]
=
\sum_x
xp_X(x).
}
\]

For a continuous random variable,

\[
\boxed{
\mathbb{E}[X]
=
\int_{-\infty}^{\infty}
xf_X(x)\,dx.
}
\]

The expectation of a function is obtained through

\[
\boxed{
\mathbb{E}[g(X)].
}
\]

Linearity gives

\[
\boxed{
\mathbb{E}[aX+bY+c]
=
a\mathbb{E}[X]
+
b\mathbb{E}[Y]
+
c.
}
\]

For an indicator,

\[
\boxed{
\mathbb{E}[\mathbf{1}_A]
=
P(A).
}
\]

Variance is

\[
\boxed{
\operatorname{Var}(X)
=
\mathbb{E}[(X-\mu)^2].
}
\]

Equivalently,

\[
\boxed{
\operatorname{Var}(X)
=
\mathbb{E}[X^2]-\mu^2.
}
\]

Linear transformations satisfy

\[
\boxed{
\mathbb{E}[aX+b]
=
a\mathbb{E}[X]+b
}
\]

and

\[
\boxed{
\operatorname{Var}(aX+b)
=
a^2\operatorname{Var}(X).
}
\]

Covariance is

\[
\boxed{
\operatorname{Cov}(X,Y)
=
\mathbb{E}[XY]
-
\mathbb{E}[X]\mathbb{E}[Y].
}
\]

Variance of a sum is

\[
\boxed{
\operatorname{Var}(X+Y)
=
\operatorname{Var}(X)
+
\operatorname{Var}(Y)
+
2\operatorname{Cov}(X,Y).
}
\]

Correlation is

\[
\boxed{
\rho_{X,Y}
=
\frac{
\operatorname{Cov}(X,Y)
}{
\sigma_X\sigma_Y
}.
}
\]

It satisfies

\[
\boxed{
-1\leq\rho_{X,Y}\leq1.
}
\]

Conditional expectation satisfies

\[
\boxed{
\mathbb{E}[X]
=
\mathbb{E}
[
\mathbb{E}[X\mid Y]
].
}
\]

Total variance decomposes as

\[
\boxed{
\operatorname{Var}(X)
=
\mathbb{E}
[
\operatorname{Var}(X\mid Y)
]
+
\operatorname{Var}
(
\mathbb{E}[X\mid Y]
).
}
\]

Markov's inequality gives

\[
\boxed{
P(X\geq a)
\leq
\frac{\mathbb{E}[X]}{a}
}
\]

for nonnegative \(X\).

Chebyshev's inequality gives

\[
\boxed{
P(|X-\mu|\geq a)
\leq
\frac{\sigma^2}{a^2}.
}
\]

Moments therefore provide a systematic language for

\[
\boxed{
\text{location}
+
\text{spread}
+
\text{shape}
+
\text{dependence}
+
\text{probability bounds}.
}
\]

%%%%%%%%%%%%%%%%%%%%%%%%%%%%%%%%%%%%%%%%%%%%%%%%%%%%%%%%%%%%
\subsection{Section Connections}
\label{subsec:expectation-moments-connections}
%%%%%%%%%%%%%%%%%%%%%%%%%%%%%%%%%%%%%%%%%%%%%%%%%%%%%%%%%%%%

\chapterconnection{
Expectation and Moments converts the distributions developed in
Sections~\ref{sec:discrete-distributions},
\ref{sec:continuous-distributions}, and
\ref{sec:joint-distributions} into numerical summaries of location,
spread, shape, and dependence. The next section, Functions of Random
Variables, develops systematic transformation methods for constructing
new distributions from existing random variables. Probability
Transforms then encodes moments and distributions through generating
functions, moment-generating functions, characteristic functions, and
related transforms. The Laws of Large Numbers explain why sample
averages approach expected values, while Central Limit Theorems explain
the asymptotic behavior of centered and standardized sums. Statistics
uses expectation, variance, covariance, and conditional expectation
throughout estimation, regression, experimental design, and inference.
}

%%%%%%%%%%%%%%%%%%%%%%%%%%%%%%%%%%%%%%%%%%%%%%%%%%%%%%%%%%%%
% SECTION SOURCE TRACKING
%%%%%%%%%%%%%%%%%%%%%%%%%%%%%%%%%%%%%%%%%%%%%%%%%%%%%%%%%%%%

%------------------------------------------------------------
% 7.6 Expectation and Moments
%------------------------------------------------------------
%
% Source basis:
%   - Standard probability theory.
%   - Standard expectation and moment theory.
%   - Standard covariance and conditional-expectation theory.
%
% Used for:
%   - discrete expectation;
%   - continuous expectation;
%   - expectation existence;
%   - law of the unconscious statistician;
%   - linearity of expectation;
%   - indicator expectations;
%   - counting with indicators;
%   - tail-sum formulas;
%   - raw moments;
%   - central moments;
%   - variance;
%   - standard deviation;
%   - transformations of mean and variance;
%   - means and variances of Bernoulli distributions;
%   - binomial moments;
%   - geometric moments;
%   - negative binomial moments;
%   - hypergeometric moments;
%   - Poisson moments;
%   - uniform moments;
%   - exponential moments;
%   - gamma moments;
%   - normal moments;
%   - beta moments;
%   - chi-square moments;
%   - covariance;
%   - covariance identities;
%   - variance of sums;
%   - independence and covariance;
%   - correlation;
%   - covariance matrices;
%   - joint expectations;
%   - conditional expectation;
%   - law of total expectation;
%   - law of total variance;
%   - higher moments;
%   - skewness;
%   - kurtosis;
%   - factorial moments;
%   - moment existence;
%   - Cauchy and Pareto moment behavior;
%   - Markov's inequality;
%   - Chebyshev's inequality;
%   - Jensen's inequality preview;
%   - Cauchy--Schwarz;
%   - standardization;
%   - PGF and MGF moment connections;
%   - cumulant preview;
%   - applications and exercises.
%
% Notes:
%   - Transformations of random variables are developed fully in
%     Section 7.7.
%   - Probability transforms and cumulants are developed in Section 7.8.
%   - Laws of Large Numbers later formalize long-run convergence to
%     expected values.
%   - Central Limit Theorems later explain standardized sums and sample
%     means.
%   - Conditional expectation receives its full measure-theoretic
%     treatment in Section 7.11.
%
%%%%%%%%%%%%%%%%%%%%%%%%%%%%%%%%%%%%%%%%%%%%%%%%%%%%%%%%%%%%